% !TeX program = xelatex
% 数学 LaTeX 排版 Skill v4.0.0；版式保持 v3 金标准. XeLaTeX 编译至少两次. 不分发字体.
% WZ-LOCKED-CLASS-BEGIN
\begin{filecontents*}[overwrite]{elegantbook-original-adapter.cls}
\NeedsTeXFormat{LaTeX2e}
\ProvidesClass{elegantbook-original-adapter}[2026/09/23 v3.0 fixed mathematical book environments]
\LoadClass[10pt,letterpaper,oneside]{book}
\RequirePackage[UTF8,scheme=plain,fontset=fandol]{ctex}
\RequirePackage{fontspec}
\setmainfont{cmunrm.otf}[BoldFont=cmunbx.otf,ItalicFont=cmunti.otf,BoldItalicFont=cmunbi.otf]
\setCJKfamilyfont{sourceheader}[ItalicFont=FandolKai-Regular.otf]{FandolKai-Regular.otf}
\RequirePackage{amsmath,amssymb,mathrsfs}
\RequirePackage{amsthm,etoolbox}
\RequirePackage{xcolor,geometry,setspace,fancyhdr,enumitem,needspace,xurl}
\definecolor{structurecolor}{HTML}{880E4F}
\definecolor{main}{HTML}{9C27B0}
\geometry{paperwidth=612bp,paperheight=792bp,left=20mm,right=20mm,top=95.04bp,bottom=20mm,headheight=12pt,headsep=25pt,footskip=30pt}
\setstretch{1.6}
\setlength{\parindent}{2em}
\setlength{\parskip}{0pt}
\setlist{nosep,leftmargin=2em}
\AtBeginDocument{\setlength{\abovedisplayskip}{4pt plus 1pt minus 1pt}\setlength{\belowdisplayskip}{4pt plus 1pt minus 1pt}\setlength{\abovedisplayshortskip}{2pt}\setlength{\belowdisplayshortskip}{3pt}}
\fancyhf{}
\fancyhead[L]{{\itshape\CJKfamily{sourceheader}\rightmark}}
\fancyhead[R]{{\itshape\CJKfamily{sourceheader}\leftmark}}
\fancyfoot[C]{\thepage}
\renewcommand{\headrulewidth}{0.4pt}
\renewcommand{\headrule}{{\color{structurecolor}\hrule height\headrulewidth width\headwidth\vskip-\headrulewidth}}
\pagestyle{fancy}
\fancypagestyle{cover}{\fancyhf{}\fancyfoot[C]{\thepage}\renewcommand{\headrulewidth}{0pt}\renewcommand{\headrule}{}}
\RequirePackage{hyperref}
\hypersetup{unicode=true,colorlinks=true,linkcolor=structurecolor,urlcolor=structurecolor,citecolor=structurecolor,linktoc=section,bookmarksnumbered=true,bookmarksopen=true,pdfborder={0 0 0},pdfstartview=FitH}
\setcounter{tocdepth}{1}
\renewcommand*\l@section{\@dottedtocline{1}{1.5em}{2.4em}}
\renewcommand{\thesection}{\thechapter.\arabic{section}}
\newcommand{\originalchapter}[1]{\par\addvspace{14pt}\Needspace{5\baselineskip}\refstepcounter{chapter}\setcounter{section}{0}\addcontentsline{toc}{chapter}{\protect\numberline{\thechapter}#1}\markboth{\thechapter\quad #1}{}\begingroup\centering\color{structurecolor}\bfseries\fontsize{14.4}{20}\selectfont\thechapter\quad #1\par\endgroup\nobreak\vspace{8pt}}
\newcommand{\originalsection}[1]{\par\addvspace{9pt}\Needspace{4\baselineskip}\refstepcounter{section}\addcontentsline{toc}{section}{\protect\hspace*{1.5em}\protect\numberline{\protect\textcolor{structurecolor}{\thesection}}#1}\markright{\thesection\quad #1}\noindent{\fontsize{12}{17}\selectfont\color{structurecolor}\thesection\quad}{\bfseries\fontsize{12}{17}\selectfont #1}\par\nobreak\vspace{4pt}}

% Semantic environments are registered once here, not re-designed in manuscripts.
\newtheoremstyle{wzstatement}{6pt}{5pt}
  {\normalfont\kaishu}{0pt}{\normalfont\bfseries\color{structurecolor}}
  {}{.7em}{\thmname{#1}\thmnumber{ #2}\thmnote{\normalfont\ (#3)}}
\newtheoremstyle{wzdefinition}{6pt}{5pt}
  {\normalfont\songti}{0pt}{\normalfont\bfseries\color{structurecolor}}
  {}{.7em}{\thmname{#1}\thmnumber{ #2}\thmnote{\normalfont\ (#3)}}
\newtheoremstyle{wzproblem}{7pt}{5pt}
  {\normalfont\songti}{2em}{\normalfont\bfseries\color{main}}
  {}{.7em}{\thmname{#1}\thmnumber{ #2}}
\newtheoremstyle{wzremark}{4pt}{4pt}
  {\normalfont\songti}{0pt}{\normalfont\bfseries}
  {}{.7em}{\thmname{#1}}
\theoremstyle{wzstatement}
\newtheorem{theorem}{定理}[chapter]
\newtheorem{lemma}[theorem]{引理}
\newtheorem{proposition}[theorem]{命题}
\newtheorem{corollary}[theorem]{推论}
\theoremstyle{wzdefinition}
\newtheorem{definition}[theorem]{定义}
\theoremstyle{wzproblem}
\newtheorem{example}{例题}[chapter]
\newtheorem{exercise}{习题}[chapter]
\theoremstyle{wzremark}
\newtheorem*{remark}{注}
\renewcommand{\proofname}{证明}
% Retain the native amsthm QED stack; only adjust the fixed typography.
\expandafter\patchcmd\csname\string\proof\endcsname{\normalfont \topsep6\p@\@plus6\p@\relax}
  {\normalfont\songti \topsep5\p@\@plus1\p@\relax}
  {}{\ClassError{elegantbook-original-adapter}{Cannot patch proof spacing}{Check the amsthm version.}}
\expandafter\patchcmd\csname\string\proof\endcsname{\itshape}{\normalfont\bfseries}
  {}{\ClassError{elegantbook-original-adapter}{Cannot patch proof heading}{Check the amsthm version.}}
\expandafter\patchcmd\csname\string\proof\endcsname{\ignorespaces}
  {\setlength{\parindent}{2em}\ignorespaces}
  {}{\ClassError{elegantbook-original-adapter}{Cannot patch proof paragraphs}{Check the amsthm version.}}
\AtBeginEnvironment{proof}{\par\Needspace{3\baselineskip}}
\renewcommand{\qedsymbol}{\ensuremath{\square}}
\newenvironment{solution}{\begin{proof}[解答]}{\end{proof}}
\newcommand{\wzcheckchapter}{\ifnum\value{chapter}<1
  \ClassError{elegantbook-original-adapter}{Numbered mathematics before chapter 1}{Move it after the first originalchapter.}\fi}

\newcommand{\wzregisterguard}[1]{\AtBeginEnvironment{#1}{\wzcheckchapter\par\Needspace{3\baselineskip}}}
\forcsvlist{\wzregisterguard}{theorem,lemma,proposition,corollary,definition,example,exercise}
% Front matter and bibliography are not numbered mathematical chapters.
\newcommand{\originalpreface}{\par\addvspace{10pt}\Needspace{5\baselineskip}
  \noindent{\bfseries\fontsize{12}{17}\selectfont 前言}\par\nobreak\vspace{4pt}}
\newcommand{\originalcontents}{\par\addvspace{14pt}
  \begin{center}{\color{structurecolor}\bfseries\fontsize{14.4}{20}\selectfont 目录}\end{center}
  \vspace{-10pt}\@starttoc{toc}}
\renewcommand{\bibname}{参考文献}
\renewenvironment{thebibliography}[1]{
  \par\addvspace{12pt}\Needspace{3\baselineskip}\phantomsection
  \addcontentsline{toc}{chapter}{\bibname}
  \markboth{\bibname}{}
  \noindent{\color{structurecolor}\bfseries\fontsize{14.4}{20}\selectfont\bibname}\par\nobreak\vspace{6pt}
  \list{\@biblabel{\@arabic\c@enumiv}}
    {\settowidth\labelwidth{\@biblabel{#1}}\leftmargin\labelwidth
     \advance\leftmargin\labelsep\usecounter{enumiv}
     \let\p@enumiv\@empty\renewcommand\theenumiv{\@arabic\c@enumiv}
     \setlength{\itemsep}{3pt}\setlength{\parsep}{0pt}}
  \sloppy\clubpenalty4000\widowpenalty4000\sfcode`\.=1000\relax
}{\def\@noitemerr{\@latex@warning{Empty bibliography}}\endlist}

\numberwithin{equation}{chapter}
\renewcommand{\theequation}{\thechapter.\arabic{equation}}
\allowdisplaybreaks[2]
\emergencystretch=2em
\clubpenalty=6000
\widowpenalty=6000
\raggedbottom
\endinput
\end{filecontents*}
% WZ-LOCKED-CLASS-END
\documentclass{elegantbook-original-adapter}
% ===== 元数据内容槽 =====
\newcommand{\BookTitle}{曲线与曲面的微分几何}
\newcommand{\BookAuthor}{Paul Seidel 原课程讲义; 中文重构与增补}
\newcommand{\BookDate}{MIT 18.950, Fall 2008}
\hypersetup{pdftitle={曲线与曲面的微分几何: MIT 18.950 中文重构本},pdfauthor={Paul Seidel; 中文整理}}
\usepackage{tikz,pgfplots,booktabs,longtable}
\pgfplotsset{compat=1.18}
\usetikzlibrary{arrows.meta,calc,positioning}
\newcommand{\R}{\mathbb R}
\newcommand{\Z}{\mathbb Z}
\newcommand{\C}{\mathbb C}
\newcommand{\dd}{\,\mathrm d}
\newcommand{\eps}{\varepsilon}
\newcommand{\nablaS}{\nabla}
\newcommand{\II}{\mathrm{II}}
\newcommand{\I}{\mathrm{I}}
\newcommand{\tr}{\operatorname{tr}}
\newcommand{\grad}{\operatorname{grad}}
\newcommand{\dist}{\operatorname{dist}}
\newcommand{\id}{\operatorname{id}}
\newcommand{\Area}{\operatorname{Area}}
\newcommand{\Exp}{\operatorname{Exp}}
\newcommand{\sech}{\operatorname{sech}}
\newcommand{\arctanh}{\operatorname{arctanh}}
\begin{document}
\frontmatter
\begin{center}
{\color{structurecolor}\bfseries\fontsize{18}{25}\selectfont\BookTitle\par}
{\bfseries MIT 18.950 中文重构本\par}
\BookAuthor\par
\BookDate\par
\end{center}
\originalpreface
曲线的弯曲由切向量的变化衡量, 曲面的弯曲由法向量的变化衡量. 把这两种变化同长度、角度和面积联系起来, 才能区分物体在空间中的形状与物体本身的内蕴几何. 本书以本科平面曲线、空间曲线和三维欧氏空间中的光滑曲面为主线, 从弧长与 Frenet 标架进入基本形式、Gauss 曲率、测地线和 Gauss--Bonnet 定理.

读者需会多元微积分、线性代数和初等常微分方程. 逆函数与隐函数定理、带参数光滑常微分方程的局部存在唯一性、平面 Green 公式作为分析先修使用; 需要的应用条件在正文列明. 曲面的有限光滑三角剖分、Euler 示性数与球面的单连通性作为拓扑工具引用, 不在本册重写一般拓扑、基本群和覆盖理论. 这些工具出现前会说明作用. 曲率公式、Frenet 方程、Gauss 方程、测地线变分与局部极小性、局部及整体 Gauss--Bonnet 的几何推导均在正文展开.

主要依据为 Paul Seidel 的 \href{https://ocw.mit.edu/courses/18-950-differential-geometry-fall-2008/pages/lecture-notes/}{MIT OpenCourseWare 18.950 Differential Geometry, Fall 2008} 四份官方公开讲义, 共60个物理页, 对应第1--41讲. 官方范围包括平面曲线、任意维超曲面、整体几何与长度距离, 也包含高维 Frenet 理论、Minkowski 几何、次数和度量曲率. 本书选取其中的曲线与二维曲面主干, 对第5--6讲作三维特化, 并补充细节、经典算例与独立练习; 不是四份原件的逐页直译或全部课程内容的替代品. 原文的无证明或概略证明结论, 在本书主干范围内补成证明, 或明确列为所用先修工具. 高维外代数曲率、Lorentz 几何、Whitney 公式、一般次数理论、CAT 与 Busemann 空间不在本册范围. 双曲平面以二维内蕴度量介绍.

原课程把平均曲率定义为主曲率之和. 本书采用曲面教材常用的 $H=(k_1+k_2)/2$, 因而本书 $2H$ 对应原讲义的 mean curvature. 统一取形算子 $S=-\dd N$, 第二基本形式 $\II(V,W)=\langle SV,W\rangle$; 外法向球面的主曲率为负. 平面有向曲率、空间曲率与测地曲率分别定义, 不混用符号的正负约定. 所有图均按本书公式用 TikZ/PGFPlots 独立绘制, 不复用课程首页外部图片.

原讲义作者为 Paul Seidel, 原机构为 Massachusetts Institute of Technology, 年份为2008. 官方文件保留其 MIT OCW 来源页. MIT OCW 当前默认许可为 \href{https://creativecommons.org/licenses/by-nc-sa/4.0/}{CC BY-NC-SA 4.0 International}; 本中文改编按相同许可用于非商业分享. 翻译、章节重组、证明增补、符号约定和另设练习为本整理本的改动, 不表示 MIT 或原作者认可本项目. 单独权利声明优先; 第三方受限材料只留官方链接. \href{https://ocw.mit.edu/pages/privacy-and-terms-of-use/}{官方使用条款}、逐文件校验和与权利核查另见来源清单. 本版核对日期为2026-10-08.
\clearpage
\originalcontents
\clearpage
\mainmatter
\originalchapter{弧长与平面曲率}
\originalsection{正则曲线与弧长}
曲线不只是一个点集. 同一个圆可以匀速或变速走过, 也可以改变行走方向. 几何量应能区分形状与参数造成的差异. 本书用 $\langle\cdot,\cdot\rangle$ 表示欧氏内积, $|\cdot|$ 表示相应范数.
\begin{definition}
设 $I\subset\R$ 为区间. 给定整数 $r\ge1$, 一个 $C^r$ 映射 $\gamma:I\to\R^n$ 若处处满足 $\gamma'(t)\ne0$, 称为正则曲线. 速度为 $v(t)=|\gamma'(t)|$. 区间 $[a,b]\subset I$ 上的长度为 $L(\gamma)=\int_a^b v(t)\dd t$.
\end{definition}
后文默认对象光滑, 即 $C^\infty$; 单项结论只需有限阶导数时会指出. 例如曲率需要 $C^2$, 挠率需要 $C^3$. 正则性使曲线具有一个确定的切线方向, 并允许弧长成为新参数.
\begin{proposition}\label{prop:arclength}
取 $t_0\in I$, 定义 $s(t)=\int_{t_0}^t v(u)\dd u$. 则 $s$ 严格递增, 映射到区间 $J=s(I)$, 逆函数 $t=t(s)$ 与原曲线同阶光滑. 曲线 $\alpha(s)=\gamma(t(s))$ 满足 $|\alpha'(s)|=1$.
\end{proposition}
\begin{proof}
有 $s'(t)=v(t)>0$, 故 $s$ 严格递增而且单射. 逆函数定理在每一点给出局部光滑逆函数, 这些逆函数由单射性在交集上一致, 因而拼成 $J$ 上的逆. 链式法则给出 $\alpha'=\gamma'/v$, 所以其长度为1. 对有限闭区间, 端点处的光滑性理解为来自邻近开区间的限制.
\end{proof}
若 $\widetilde\gamma=\gamma\circ\phi$, $\phi$ 是区间间光滑微分同胚, 则 $|\widetilde\gamma'|=v(\phi)|\phi'|$. 换元公式说明长度不变. $\phi'>0$ 保持定向, $\phi'<0$ 反转定向. 长度只依赖所走弧段及其重数; 重复绕一圈仍会增加长度.
\begin{example}
比较 $\gamma(t)=(t^3,0)$ 与 $\eta(t)=(t,t^2)$ 在 $t=0$ 附近的参数性质.
\end{example}
\begin{solution}
$\gamma'(0)=0$, 所以它不是所给参数下的正则曲线, 尽管其像是一条直线. 改用横坐标能正则描述同一点集, 但 $t\mapsto t^3$ 在0的逆不是光滑微分同胚. 对 $\eta$, 有 $\eta'=(1,2t)$, 处处非零. 弧长 $s(t)=\int_0^t\sqrt{1+4u^2}\dd u$ 严格递增; 无须写出逆的初等表达式, 已能合法地使用弧长参数.
\end{solution}
\originalsection{有向曲率与 Frenet 方程}
在定向平面中记 $J(x,y)=(-y,x)$, 即逆时针旋转直角. 若 $T$ 是单位切向量, 左法向量为 $n=JT$. 由 $\langle T,T\rangle=1$ 得 $T'\perp T$, 所以切向量的变化只有法向分量.
\begin{definition}
单位速度平面曲线的有向曲率为 $k(s)=\det(T(s),T'(s))$. 对任意正则参数, 定义
\begin{equation}\label{eq:plane-curvature}
k(t)=\frac{\det(\gamma'(t),\gamma''(t))}{|\gamma'(t)|^3}.
\end{equation}
\end{definition}
\begin{proposition}\label{prop:plane-frenet}
对任意参数, 有 $\dd T/\dd t=v k n$, $\dd n/\dd t=-v k T$. 对弧长参数, 有 $T'=k n$, $n'=-kT$.
\end{proposition}
\begin{proof}
写 $\gamma'=vT$, 则 $\gamma''=v'T+vT'$. 由于 $T'$ 与 $T$ 垂直, 可写成 $aJT$. 取行列式得 $\det(\gamma',\gamma'')=v^2a$, 所以 $a=vk$. 再用 $J^2=-\id$ 得第二式. 在弧长参数下 $v=1$.
\end{proof}
曲率的量纲为长度的倒数. 正值表示切向量向左转, 负值表示向右转; $|k|$ 衡量弯曲强度. 换参时二阶导数中的 $\phi''\gamma'$ 不贡献行列式, 从而
\[
 k_{\widetilde\gamma}(u)=\operatorname{sgn}(\phi'(u))k_\gamma(\phi(u)).
\]
故保持定向时曲率不变, 反向时有向曲率变号. 平面反射同样使有向曲率变号, 平移与正向旋转保持它.
\begin{theorem}\label{thm:plane-fundamental}
给定区间 $J$ 上的连续函数 $k$、点 $p$ 和单位向量 $T_0$, 存在唯一 $C^2$ 单位速度平面曲线 $\alpha$ 具有该曲率并满足 $\alpha(s_0)=p$, $\alpha'(s_0)=T_0$. 因而曲率函数确定曲线, 精确到平移和正向旋转.
\end{theorem}
\begin{proof}
取 $T_0=(\cos\theta_0,\sin\theta_0)$, 定义
\[
 \theta(s)=\theta_0+\int_{s_0}^s k(u)\dd u,\qquad
 \alpha(s)=p+\int_{s_0}^s(\cos\theta(u),\sin\theta(u))\dd u.
\]
直接微分得单位速度与 $T'=kJT$. 为证明唯一性, 任意另一单位切向量 $\widetilde T$ 满足同一线性方程 $\widetilde T'=kJ\widetilde T$ 和相同初值. 差向量 $D$ 满足 $D'=kJD$, 故 $(|D|^2)'=2k\langle D,JD\rangle=0$. 由初值为零可知切向量一致, 积分后曲线一致.
\end{proof}
\begin{corollary}
曲率恒为0的连通正则平面曲线位于直线上. 曲率恒为非零常数 $c$ 的曲线位于半径 $1/|c|$ 的圆上.
\end{corollary}
\begin{proof}
弧长化后, 前一种情形 $T'=0$. 后一种情形 $\alpha+n/c$ 的导数为 $T-T=0$, 故圆心固定, 且 $|\alpha-\text{圆心}|=1/|c|$.
\end{proof}
\originalsection{密切圆与计算}
在点 $s_0$ 附近, 单位速度曲线的 Taylor 展开为
\[
\alpha(s_0+h)=\alpha(s_0)+hT(s_0)+\tfrac12h^2 k(s_0)n(s_0)+o(h^2).
\]
一次项给切线, 二次项给曲率. 当 $k(s_0)\ne0$ 时, 与这两个导数同时相同的圆称为密切圆, 圆心为 $\alpha(s_0)+n(s_0)/k(s_0)$, 半径为 $1/|k(s_0)|$. 唯一性来自圆心必须位于法线上, 且圆的曲率固定其有向距离. $k=0$ 时密切对象是切线, 无有限半径的密切圆.
\begin{center}
\begin{tikzpicture}[scale=1.15,>=Latex]
\draw[->] (-2.3,0)--(2.3,0) node[right] {$x$};
\draw[->] (0,-.3)--(0,2.35) node[above] {$y$};
\draw[thick,structurecolor,domain=-1.35:1.35,samples=60] plot (\x,{\x*\x/2});
\draw[dashed] (0,1) circle (1);
\draw[->,thick] (0,0)--(1.0,0) node[below] {$T$};
\draw[->,thick] (0,0)--(0,.65) node[left] {$n$};
\fill (0,1) circle (1.4pt) node[right=4pt] {圆心};
\node at (1.7,1.9) {$y=x^2/2$};
\node at (-1.65,1.8) {密切圆};
\end{tikzpicture}
\end{center}
图中抛物线在原点有 $k=1$. 图示只表达二阶接触, 不表示曲线与圆在邻域重合.
\begin{example}
求椭圆 $\gamma(t)=(a\cos t,b\sin t)$, $a\ge b>0$, 的曲率和曲率极值.
\end{example}
\begin{solution}
有 $v^2=a^2\sin^2t+b^2\cos^2t$, 而 $\det(\gamma',\gamma'')=ab$. 所以
\[
 k(t)=\frac{ab}{(a^2\sin^2t+b^2\cos^2t)^{3/2}}.
\]
分母最小时曲率最大: 在 $(\pm a,0)$ 有 $k=a/b^2$; 在 $(0,\pm b)$ 有 $k=b/a^2$. 若 $a>b$, 这四点是全部曲率极值点. 不能把参数的二阶导数范数直接当曲率, 因为此参数一般不是单位速度.
\end{solution}
图像曲线 $\gamma(x)=(x,f(x))$ 的曲率为 $f''/(1+f'^2)^{3/2}$. 对隐式曲线也有不必先解出图像的公式.
\begin{proposition}
设 $F\in C^2(U)$, $\nabla F\ne0$ 在正则水平曲线 $F=c$ 上. 选择单位切向量 $T=J\nabla F/|\nabla F|$, 则
\[
 k=\frac{\operatorname{Hess}F(J\nabla F,J\nabla F)}{|\nabla F|^3}.
\]
\end{proposition}
\begin{proof}
弧长微分 $F(\alpha(s))=c$ 两次, 得 $\operatorname{Hess}F(T,T)+\langle\nabla F,\alpha''\rangle=0$. 所选左法向量为 $n=-\nabla F/|\nabla F|$, 代入 $\alpha''=kn$ 即得公式. 反转定向会使结果变号.
\end{proof}
\originalsection{习题}
\begin{exercise}\label{ex:1a}
求 $\gamma(t)=(e^t\cos t,e^t\sin t)$ 的速度与有向曲率, 并求 $t=0$ 处密切圆的圆心和半径.
\end{exercise}
\begin{exercise}\label{ex:1b}
设正则平面曲线曲率处处非零. 以弧长表示其密切圆圆心轨迹 $C(s)=\alpha(s)+n(s)/k(s)$. 求 $C'$, 并说明曲率的临界点如何影响此轨迹的正则性.
\end{exercise}
\begin{exercise}\label{ex:1c}
对水平集 $x^2+2y^2=1$, 取由 $T=J\nabla F/|\nabla F|$ 确定的方向, 计算 $k$ 在 $(1,0)$ 与 $(0,1/\sqrt2)$ 的值, 与椭圆公式核对.
\end{exercise}

\originalchapter{闭曲线的整体几何}
\originalsection{切角与旋转数}
局部曲率告诉我们切线在一小段上怎样转动, 但闭合条件会约束整条曲线累计转了多少. 一个闭曲线在这里指周期正则映射 $\gamma:\R\to\R^2$, $\gamma(t+P)=\gamma(t)$. 若 $[0,P)$ 上不同参数不对应同一点, 称为简单闭曲线.
\begin{lemma}\label{lem:angle}
设 $T:I\to S^1$ 为 $C^1$ 映射. 可取 $C^1$ 实函数 $\theta$ 使 $T=(\cos\theta,\sin\theta)$; 指定 $\theta(t_0)$ 后唯一, 且 $\theta'=\det(T,T')$.
\end{lemma}
\begin{proof}
选择一个与 $T(t_0)$ 对应的角 $\theta_0$, 定义 $\theta(t)=\theta_0+\int_{t_0}^t\det(T,T')\dd u$. 由单位向量的微分分解, 原 $T$ 与由该角定义的向量都满足 $V'=\det(T,T')JV$. 前章的线性方程唯一性说明它们相同. 两个角函数的差连续且属于 $2\pi\Z$, 在区间上必为常数.
\end{proof}
\begin{definition}
闭曲线的旋转数为 $m=(\theta(P)-\theta(0))/(2\pi)$, 其中 $\theta$ 为单位切向量的连续切角. 总有向曲率为 $\int_0^P k(t)v(t)\dd t$.
\end{definition}
由于 $T(P)=T(0)$, 旋转数是整数. 引理给出
\begin{equation}\label{eq:turning}
 \int_0^P k(t)v(t)\dd t=2\pi m.
\end{equation}
闭合本身并不强迫 $m=1$. 圆重复走 $q$ 圈有 $m=q$, 而有自交的闭曲线可以有 $m=0$. 绕某个不在曲线上的点的绕数, 则是位置向量 $\gamma-p$ 的角变化; 它与切向量旋转数是两个不同量.
\begin{example}
求 $\gamma(t)=(\sin t,\sin 2t)$, $0\le t\le2\pi$, 的旋转数.
\end{example}
\begin{solution}
$\gamma'=(\cos t,2\cos2t)$ 不会同时为零, 故曲线正则. 参数平移 $t\mapsto t+\pi$ 时, 速度不变, 而
\[
 \det(\gamma',\gamma'')=-2\sin t\,(1+2\cos^2t)
\]
变号. 因而 $kv=\det(\gamma',\gamma'')/v^2$ 在两个半周期的积分相消, 总曲率为0, 旋转数为0. 它在原点自交, 所以不属于下面的简单闭曲线结论.
\end{solution}
\begin{center}
\begin{tikzpicture}
\begin{axis}[width=7.5cm,height=4.1cm,axis equal image,axis lines=middle,xtick=\empty,ytick=\empty,xmin=-1.25,xmax=1.25,ymin=-1.25,ymax=1.25,xlabel={$x$},ylabel={$y$},xlabel style={at={(axis description cs:1,.5)},anchor=west},ylabel style={at={(axis description cs:.5,1)},anchor=south}]
\addplot[thick,structurecolor,domain=0:360,samples=181] ({sin(x)},{sin(2*x)});
\end{axis}
\end{tikzpicture}
\end{center}
\originalsection{转角定理与凸性}
凸区域指任意两点间线段仍在区域内; 严格凸指凸区域的边界不含非平凡线段. 支撑线指经过边界点且把整个区域置于一侧的直线. 这里使用平面简单闭曲线分割内外区域的 Jordan 定理, 以及简单多边形可三角剖分的拓扑事实. 二者均不是曲率公式的推论. 曲线沿区域在左侧的方向称为正向边界.
\begin{theorem}[转角定理]\label{thm:turning}
$C^2$ 正则简单闭平面曲线若按正向边界定向, 则总有向曲率为 $2\pi$; 反向时为 $-2\pi$.
\end{theorem}
\begin{proof}
先看有 $q$ 个顶点的正向简单多边形. 三角剖分给出内部角之和 $(q-2)\pi$, 凹顶点的内角允许大于 $\pi$. 在顶点把切向量转入下一条边的有向外角为 $\pi-\alpha_j\in(-\pi,\pi)$, 所以外角总和为 $2\pi$.

将光滑曲线弧长化, 周期为 $L$. 单位切向量连续, 且其导数有界. 把周期分为足够细的小段, 用弦连接端点. 这些弦的单位方向与该小段上的切向量一致地接近: 这是 $\alpha(b)-\alpha(a)=\int_a^b T\dd s$ 及 $T$ 一致连续的结果.

还需保证细弦构成的多边形简单. 为此在 $s\in\R/L\Z$ 上考虑 $F(s,r)=\alpha(s)+rn(s)$. 有 $F_s=(1-rk)T,F_r=n$. 缩小带宽使 $1-rk>0$, 在 $r=0$ 的导数以 $T,n$ 为列向量, 可逆; 逆函数定理给出局部带状邻域. 紧致性使局部宽度可以统一缩小. 对远离圆周参数对角线的 $(s,t)$, 简单性及紧致性使 $|\alpha(s)-\alpha(t)|$ 有正下界; 再缩小宽度便排除远处的重叠. 因而 $F$ 给出整条曲线的管状带, 每一点有唯一投影参数 $s\in\R/L\Z$. 足够细的每条弦处于相应局部带中, 投影参数与原小弧接近, 其方向与该投影点处的 $T$ 的内积为正. 若弦以参数 $u$ 表示, 则 $\dd s/\dd u=\langle\text{弦速度},T(s)\rangle/(1-rk(s))>0$, 故投影参数沿弦严格增加, 并以两端点的参数为起终值. 相邻弦只在共同端点相遇, 不相邻弦的投影参数区间内部不交. 多边形因此简单, 且方向与原曲线一致.

记细分点为 $s_j$, 把每条弦的角 $\beta_j$ 选成靠近 $\theta(s_j)$ 的实数分支. 写 $\theta(L)=\theta(0)+2\pi m$, 闭合时设下一周的 $\beta_q=\beta_0+2\pi m$. 细分足够细使全部相邻角差 $\beta_{j+1}-\beta_j$ 的绝对值小于 $\pi$, 包括最后一个差; 它们正是多边形的有向外角. 相加后为 $2\pi m=\theta(L)-\theta(0)$. 多边形外角总和已为 $2\pi$, 故 $\theta(L)-\theta(0)=2\pi$. 由式\eqref{eq:turning} 得结论; 反向结论由换参变号得到.
\end{proof}
\begin{corollary}
简单闭平面曲线满足 $\int|k|\dd s\ge2\pi$. 等号成立当且仅当相应定向下 $k$ 不变号.
\end{corollary}
\begin{proof}
三角不等式给出 $\int|k|\dd s\ge|\int k\dd s|=2\pi$. 对连续 $k$, 若同时有正值和负值, 两个有正长度的区间使不等式严格. 反之不变号时等号成立.
\end{proof}
\begin{proposition}\label{prop:strict-convex}
正向简单闭曲线若 $k>0$, 则每条切线都是支撑线, 且每个法向方向恰有一个支撑点. 因而它围成严格凸区域.
\end{proposition}
\begin{proof}
切角 $\theta(s)$ 严格增加, 由转角定理一周恰增加 $2\pi$. 对任意单位向量 $a$, 函数 $f(s)=\langle\alpha(s),a\rangle$ 的导数为 $\langle T(s),a\rangle$. 一周内它恰有一个从正变负的零点及一个从负变正的零点, 分别为唯一最大和最小. 切线垂直于 $a$ 的这两个点因此是全局支撑点. 进一步, 介于高度最大与最小之间的每条垂直于 $a$ 的直线恰与边界相交两点, 因为高度函数在两段上分别严格递增和递减. 由 Jordan 内外区分, 区域在此直线上的截面恰为两交点间的区间. 因为 $a$ 任意, 任意直线与闭区域的交都是区间或空集, 这等价于区域凸. 若边界含非平凡线段, 其曲率在该段为0, 与假设矛盾.
\end{proof}
\originalsection{支撑函数与四顶点定理}
严格正曲率时可用外法向的角 $\varphi$ 作参数. 记 $n_o(\varphi)=(\cos\varphi,\sin\varphi)$, $t(\varphi)=Jn_o$, 定义支撑函数 $h(\varphi)=\langle\alpha(\varphi),n_o(\varphi)\rangle$. 不要求原点在内部; 平移只给 $h$ 增加一阶三角函数.
\begin{proposition}\label{prop:support}
有 $\alpha=hn_o+h't$, 且 $\alpha'=(h+h'')t$. 曲率半径为 $\rho=h+h''>0$, 曲率为 $k=1/\rho$.
\end{proposition}
\begin{proof}
因 $\alpha'$ 沿 $t$ 方向, 微分 $h$ 得 $h'=\langle\alpha,t\rangle$. 在正交基 $n_o,t$ 下展开 $\alpha$ 得第一式. 再用 $n_o'=t$, $t'=-n_o$ 微分, 两个法向项相消, 得第二式. 切向角随 $\varphi$ 增加的速度为1, 弧长速度为 $\rho$, 所以 $k=1/\rho$.
\end{proof}
\begin{lemma}[初阶 Sturm--Hurwitz 引理]\label{lem:sturm}
连续 $2\pi$ 周期函数 $f$ 若与 $1,\cos\varphi,\sin\varphi$ 都积分正交, 则至少有四个不同零点.
\end{lemma}
\begin{proof}
若零点无限多则已成立. 否则若无符号变化, $f$ 与常数1的积分不可能为0, 除非恒为0. 周期函数的符号变化次数为偶数; 假设少于四个零点, 则只能有两个符号变化点 $a,b$. 取一阶三角多项式 $p(\varphi)=\cos(\varphi-(a+b)/2)-\cos((b-a)/2)$, 选代表元使 $0<b-a<2\pi$. 它只在 $a,b$ 为零, 在两弧上分别取相反符号. 必要时乘 $-1$, 使 $fp\ge0$ 且在非空开区间严格为正. 于是 $\int fp>0$, 但正交条件给出该积分为0, 矛盾.
\end{proof}
\begin{theorem}[严格正曲率的四顶点定理]
光滑简单闭平面曲线若 $k>0$, 则至少有四个不同点满足 $\dd k/\dd s=0$.
\end{theorem}
\begin{proof}
取上述法向角参数. $\rho'=h'+h'''$ 的积分为0. 分部积分并用周期性可得它与 $\cos\varphi$ 和 $\sin\varphi$ 的积分也为0: 在这两项中三阶导数恰抵消一阶导数. 由引理, $\rho'$ 至少有四个零点. 又 $\dd k/\dd s=-\rho'/\rho^3$, $\rho>0$, 零点一一对应. 圆的曲率恒定, 所有点都是此意义下的顶点.
\end{proof}
这里证明的是原讲义第10讲的严格正曲率版本, 不把它扩大为任意闭曲线的结论.
\originalsection{习题}
\begin{exercise}\label{ex:2a}
设单位速度曲线 $\alpha(s)$ 满足 $\alpha(s+L)=\alpha(s)$, 且 $k(s)>0$, $L$ 不要求是最小周期. 若 $\int_0^L k\dd s=4\pi$, 它相对于指定周期是否为简单曲线? 给出解释和一个例子.
\end{exercise}
\begin{exercise}\label{ex:2b}
设 $h(\varphi)=a+b\cos2\varphi$, $a>3|b|>0$. 由支撑函数公式构造曲线, 求曲率及其临界点, 并证明所得闭曲线简单.
\end{exercise}
\begin{exercise}\label{ex:2c}
设正向严格正曲率简单闭曲线的支撑函数为 $h$. 证明其周长为 $\int_0^{2\pi}h\dd\varphi$, 面积为 $\frac12\int_0^{2\pi}(h^2-h'^2)\dd\varphi$.
\end{exercise}

\originalchapter{空间曲线与挠率}
\originalsection{Frenet 标架}
空间曲线除弯曲强度外, 还会离开当前的密切平面. 平面有向曲率的一个实数不足以描述这两种变化. 本章以 $C^3$ 单位速度曲线 $\alpha:J\to\R^3$ 为起点.
\begin{definition}
若 $\kappa(s)=|\alpha''(s)|>0$, 定义单位切向量 $T=\alpha'$, 主法向量 $n=\alpha''/\kappa$, 副法向量 $b=T\times n$. $(T,n,b)$ 称为 Frenet 标架; 张成 $T,n$ 的平面称为密切平面. 挠率定义为 $\tau=\langle n',b\rangle$.
\end{definition}
$\kappa>0$ 是定义主法向量的条件, 不能在拐直点仍无条件使用此标架. 平面曲线的 $\kappa$ 是前章 $|k|$; 空间挠率有正负, 空间曲率非负.
\begin{theorem}[Frenet--Serret 方程]\label{thm:space-frenet}
在 $\kappa>0$ 的区间上,
\begin{equation}\label{eq:frenet3}
 T'=\kappa n,\qquad n'=-\kappa T+\tau b,\qquad b'=-\tau n.
\end{equation}
而且 $\tau=\det(\alpha',\alpha'',\alpha''')/\kappa^2$.
\end{theorem}
\begin{proof}
第一式来自定义. 微分 $\langle n,n\rangle=1$ 得 $n'\perp n$, 微分 $\langle n,T\rangle=0$ 得 $\langle n',T\rangle=-\kappa$. $b$ 分量是 $\tau$, 得第二式. 微分 $b=T\times n$ 得 $b'=T'\times n+T\times n'=-\tau n$. 最后 $\alpha'''=\kappa'n+\kappa(-\kappa T+\tau b)$, 而 $\alpha'\times\alpha''=\kappa b$, 三重积为 $\kappa^2\tau$.
\end{proof}
\begin{proposition}\label{prop:space-param}
任意正则参数下, 若 $\gamma'\times\gamma''\ne0$, 则
\[
 \kappa=\frac{|\gamma'\times\gamma''|}{|\gamma'|^3},\qquad
 \tau=\frac{\det(\gamma',\gamma'',\gamma''')}{|\gamma'\times\gamma''|^2}.
\]
两者在参数反向时都保持数值, 在空间反射时 $\tau$ 变号而 $\kappa$ 不变.
\end{proposition}
\begin{proof}
令 $\gamma'=vT$, 则 $\gamma''=v'T+v^2\kappa n$, 故叉积为 $v^3\kappa b$. 再微分, $\gamma'''$ 的 $b$ 分量为 $v^3\kappa\tau$, 所以三重积为 $v^6\kappa^2\tau$. 两个公式成立. 对任意换参, 叉积取得 $\phi'^3$ 因子, 三重积取得 $\phi'^6$ 因子; 平方或绝对值消去符号. 空间正交变换对三重积乘以其行列式, 给出最后的断言.
\end{proof}
\begin{proposition}
在连通区间上且 $\kappa>0$ 时, $\tau\equiv0$ 当且仅当曲线位于一个固定平面中.
\end{proposition}
\begin{proof}
若 $\tau=0$, 则 $b'=0$, $b$ 固定, 且 $\langle\alpha',b\rangle=0$. 因而 $\langle\alpha,b\rangle$ 为常数, 曲线位于该平面. 反之若曲线位于法向为 $a$ 的平面, $T,n$ 都与 $a$ 垂直; 连续单位副法向量必为固定的 $a$ 或 $-a$, 所以 $b'=0$, $\tau=0$.
\end{proof}
\originalsection{曲线基本定理}
平面曲线基本定理是一个角函数的积分. 空间情形则需积分一个旋转标架. 这里使用线性常微分方程的存在唯一性: 连续系数在区间上给定初值后有唯一解; 在系数连续的每个紧子区间上解可继续, 光滑系数给光滑解.
\begin{theorem}\label{thm:space-fundamental}
给定光滑函数 $\kappa>0,\tau$ 于区间 $J$, 固定 $s_0\in J$, 在 $s_0$ 处给定初始点 $p$ 和正向正交标架 $(T_0,n_0,b_0)$, 存在唯一单位速度空间曲线具有这些曲率、挠率和初值. 因而 $(\kappa,\tau)$ 确定曲线, 精确到正向欧氏刚体运动.
\end{theorem}
\begin{proof}
以列向量组成矩阵 $Q=(T,n,b)$, 解
\[
 Q'=QA,\qquad A=\begin{pmatrix}0&-\kappa&0\\\kappa&0&-\tau\\0&\tau&0\end{pmatrix},\qquad Q(s_0)=Q_0.
\]
由于 $A^{\mathsf T}=-A$, 矩阵 $M=Q^{\mathsf T}Q$ 满足 $M'=A^{\mathsf T}M+MA$. 常矩阵 $M=I_3$ 是具有相同初值的解, 唯一性使其一直成立. $\det Q$ 连续、只取 $\pm1$ 且初值为1, 故标架始终正向. 它的列向量范数为1, 不会在有限区间爆破, 因此可在整个 $J$ 上解出.

定义 $\alpha=p+\int_{s_0}^s T\dd u$. 则 $\alpha'=T$, $\alpha''=\kappa n$, 所以主法向和副法向恰为所解标架, 挠率为 $\tau$. 任意另一满足条件的曲线产生同一标架方程和初值, 唯一性给出相同曲线. 任意两个初始正向标架由唯一正向正交矩阵联系, 再加平移得到最后结论.
\end{proof}
当 $\kappa=0$ 时不能由 $n=\alpha''/\kappa$ 定义 Frenet 标架. 这不表示曲线不存在, 只表示本定理的假设不再适用; 例如直线有许多可选法向标架, 但没有上述定义的挠率.
\originalsection{圆柱螺线与一般螺线}
\begin{example}
求 $\gamma(t)=(a\cos t,a\sin t,ct)$ 的曲率和挠率, 其中 $a>0$, $c\in\R$.
\end{example}
\begin{solution}
令 $q=\sqrt{a^2+c^2}$. 有 $|\gamma'|=q$,
\[
 \gamma'\times\gamma''=(ac\sin t,-ac\cos t,a^2),\qquad
 |\gamma'\times\gamma''|=aq,
\]
三重积为 $a^2c$, 故 $\kappa=a/q^2$, $\tau=c/q^2$. 弧长参数为 $s=qt$; 当 $c=0$ 时曲线退化为圆, 挠率为0. 切向量与竖直轴的夹角满足 $\cos\beta=c/q$, 与参数无关.
\end{solution}
\begin{center}
\begin{tikzpicture}
\begin{axis}[width=8cm,height=5cm,view={38}{22},axis lines=box,clip=false,xtick=\empty,ytick=\empty,ztick=\empty,xmin=-1.4,xmax=1.4,ymin=-1.4,ymax=1.4,zmin=0,zmax=3.8]
\addplot3[thick,structurecolor,domain=0:720,samples=160,samples y=0] ({cos(x)},{sin(x)},{x/200});
\node[anchor=north] at (axis description cs:.35,-.05) {$x$};
\node[anchor=west] at (axis description cs:.92,.02) {$y$};
\node[anchor=east] at (axis description cs:-.03,.50) {$z$};
\end{axis}
\end{tikzpicture}
\end{center}
图示取 $a=1$, $c=\pi^{-1}\cdot0.9$; 竖直比例用于显示两圈结构, 曲率数值由公式确定.
\begin{theorem}[一般螺线判据]\label{thm:lancret}
单位速度曲线在连通区间上满足 $\kappa>0$. 其单位切向量与某个固定单位向量夹角恒定, 当且仅当 $\tau/\kappa$ 为常数.
\end{theorem}
\begin{proof}
若 $\tau/\kappa=d$ 为常数, 则 $U=(dT+b)/\sqrt{1+d^2}$ 满足 $U'=(d\kappa-\tau)n/\sqrt{1+d^2}=0$, 且 $\langle T,U\rangle=d/\sqrt{1+d^2}$ 固定.

反之设 $\langle T,U\rangle=a$ 恒定. 微分并用 $\kappa>0$, 得 $\langle n,U\rangle=0$, 所以 $U=aT+c(s)b$. 再微分得 $0=(a\kappa-c\tau)n+c'b$, 故 $c$ 为常数, $a\kappa=c\tau$. $c$ 不可能为0, 否则 $a\ne0$ 而 $\kappa=0$. 因而 $\tau/\kappa=a/c$ 为常数.
\end{proof}
此结论中的“螺线”不要求曲率为常数. 若曲率和挠率都为常数且 $\kappa>0$, 由基本定理, 曲线与圆或圆柱螺线的一段全等.
\originalsection{习题}
\begin{exercise}\label{ex:3a}
求 $\gamma(t)=(t,t^2,t^3)$ 的曲率和挠率. 判断它是否在任意非空开区间上为平面曲线.
\end{exercise}
\begin{exercise}\label{ex:3b}
对单位速度空间曲线 $\kappa>0$, 证明曲线位于以 $p$ 为中心、半径 $R$ 的球面上时有 $\langle\alpha-p,n\rangle=-1/\kappa$. 若进一步 $\tau\ne0$, 表示 $\alpha-p$ 在 $n,b$ 中的分量, 并得到球面半径的恒等式.
\end{exercise}
\begin{exercise}\label{ex:3c}
设 $\tau/\kappa=d$ 恒定且非零. 证明将一般螺线正交投影到垂直于固定轴 $U=(dT+b)/\sqrt{1+d^2}$ 的平面后, 投影正则, 并求其速度及无向曲率.
\end{exercise}

\originalchapter{正则曲面与第一基本形式}
\originalsection{参数片与切平面}
在平面曲线上, 一个非零导数张成切线. 在曲面上, 两个独立的一阶导数张成切平面. 参数平面上的线性代数因而可以描述曲面的长度和角度.
\begin{definition}
设 $U\subset\R^2$ 开. 光滑映射 $X:U\to\R^3$ 若 $X_u\times X_v\ne0$, 称为正则参数片. 其在参数 $(u,v)$ 处的切平面方向为 $\operatorname{span}(X_u,X_v)$. 单位法向量为 $N=(X_u\times X_v)/|X_u\times X_v|$.
\end{definition}
仅有秩条件不保证全局单射. 为描述一个作为点集的曲面, 还需要求每一点邻域中存在与参数域同胚的单射参数片. 满足此条件的 $M\subset\R^3$ 称为正则嵌入曲面; 本书整体结论默认曲面为此意义, 局部公式也适用于浸入参数片.
\begin{proposition}\label{prop:implicit}
局部点集 $F^{-1}(0)$ 若 $\nabla F\ne0$, 则为正则曲面. 反之, 每个正则嵌入曲面在一点附近可经空间坐标旋转写成图像, 因而为这样的正则水平集.
\end{proposition}
\begin{proof}
若在点处 $F_z\ne0$, 隐函数定理给出 $z=f(x,y)$, 参数片 $X(u,v)=(u,v,f(u,v))$ 的叉积为 $(-f_u,-f_v,1)$, 非零. 若别的偏导非零则交换坐标.

反之, 在参数片某点, $DX$ 的秩为2, 存在一个非零的 $2\times2$ 子式. 将空间轴旋转或交换, 可假设到 $xy$ 平面的投影 $\pi\circ X$ 导数可逆. 逆函数定理给出局部微分同胚, 换参后前两个分量就是 $(x,y)$, 第三个是 $f(x,y)$. 定义 $F(x,y,z)=z-f(x,y)$, 其梯度非零且局部零集恰为曲面.
\end{proof}
\begin{proposition}
曲面上一点 $p$ 的所有光滑曲线速度组成二维线性空间 $T_pM$. 若 $p=X(u,v)$, 则 $T_pM=\operatorname{span}(X_u,X_v)$; 若曲面由 $F=0$ 定义, 则 $T_pM=\ker\dd F_p$.
\end{proposition}
\begin{proof}
曲面曲线在局部写成 $X(u(t),v(t))$, 其速度为 $u'X_u+v'X_v$. 任意此类线性组合都可由参数直线在短时间内实现, 故等式成立. 对水平集, 微分 $F(\gamma)=0$ 得速度在 $\ker\dd F_p$ 中; 两者维数均为2, 因而相等.
\end{proof}
\begin{example}
说明球面的极坐标参数在极点失效为何不表示球面有奇点; 比较圆锥顶点.
\end{example}
\begin{solution}
半径 $R$ 的球面有 $X(\theta,\varphi)=R(\sin\theta\cos\varphi,\sin\theta\sin\varphi,\cos\theta)$. 在 $\theta=0,\pi$, $X_\varphi=0$, 此参数失去秩. 然而球面是 $x^2+y^2+z^2-R^2=0$ 的零集, 梯度在球面非零, 所以换用图像参数就恢复正则性. 圆锥 $z^2=x^2+y^2$ 在顶点则不可能有唯一切平面: 经过顶点的母线方向张成三维空间. 若它是正则曲面, 所有这些速度应落在二维切空间, 矛盾.
\end{solution}
\originalsection{长度与角度}
\begin{definition}
参数片的第一基本形式为切空间内积的拉回:
\[
 \I=a\,\dd u^2+2b\,\dd u\dd v+c\,\dd v^2,\qquad
 g=\begin{pmatrix}a&b\\b&c\end{pmatrix},
\]
其中 $a=\langle X_u,X_u\rangle$, $b=\langle X_u,X_v\rangle$, $c=\langle X_v,X_v\rangle$. 常用记号也为 $E=a,F=b,G=c$. 后文采用 $E,F,G$ 写具体算例, 用 $g_{ij}$ 写一般指标公式.
\end{definition}
对非零参数向量 $w$, 有 $w^{\mathsf T}gw=|DX(w)|^2>0$, 故 $g$ 正定. 特别 $E>0$, $G>0$, $EG-F^2>0$. 曲线 $\gamma(t)=X(u(t),v(t))$ 的长度为
\[
 L=\int\sqrt{E u'^2+2F u'v'+Gv'^2}\dd t.
\]
对两个参数方向 $w,z$, 对应空间切向量的夹角满足 $\cos\beta=(w^{\mathsf T}gz)/\sqrt{(w^{\mathsf T}gw)(z^{\mathsf T}gz)}$. 仅把参数域当欧氏平面测量会忽略曲面的度量.
\begin{proposition}\label{prop:metric-transform}
若 $\widetilde X=X\circ\phi$, $\phi$ 为局部微分同胚, 其 Jacobian 为 $A$, 则 $\widetilde g=A^{\mathsf T}(g\circ\phi)A$. 因而长度和角度不依赖所选坐标.
\end{proposition}
\begin{proof}
由链式法则 $D\widetilde X=(DX\circ\phi)A$, 取转置相乘即得矩阵公式. 一个实际切向量在两个坐标中的参数分量满足 $w=A\widetilde w$, 所以二次型数值一致.
\end{proof}
\begin{definition}
曲面间的局部微分同胚 $f:M\to\widetilde M$ 若满足 $\langle\dd f(V),\dd f(W)\rangle=\langle V,W\rangle$, 称为局部等距. 它保持曲线长度和交角. 若 $f$ 全局双射且逆光滑, 则为等距映射.
\end{definition}
圆柱可以局部展开为平面, 但整圆柱与平面不是全局双射等距. 区分局部与整体对于测地线最短性同样重要.
\originalsection{面积与定向}
\begin{proposition}\label{prop:area}
曲面面积元为 $\dd A=|X_u\times X_v|\dd u\dd v=\sqrt{EG-F^2}\dd u\dd v$. 在坐标变换下按普通面积换元规则变化, 因而定义与参数无关.
\end{proposition}
\begin{proof}
$X_u,X_v$ 张成的微小平行四边形面积是其叉积长度. 恒等式 $|X_u\times X_v|^2=|X_u|^2|X_v|^2-\langle X_u,X_v\rangle^2$ 给出根号表达式. 由命题\ref{prop:metric-transform}, $\det\widetilde g=(\det A)^2\det(g\circ\phi)$, 故 $\sqrt{\det\widetilde g}=|\det A|\sqrt{\det(g\circ\phi)}$, 与二元换元公式一致.
\end{proof}
为积分一般曲面, 可以取覆盖与分割统一把积分化成坐标积分, 对紧区域也可用有限曲面三角剖分分块积分; 块的边界面积为0, 所以不会重复计数. 不应对任意非单射参数片直接积分后称为点集的面积, 因为可能带有覆盖重数.
\begin{definition}
曲面若存在全局光滑单位法向场 $N$, 称为可定向. 选择 $N$ 即选定定向; 正向切标架 $(e_1,e_2)$ 满足 $e_1\times e_2=N$. 一个参数片与定向相容当 $X_u\times X_v$ 指向 $N$.
\end{definition}
相容坐标之间的 Jacobian 行列式为正. 反转 $N$ 不改变长度面积, 但会改变第二基本形式及有向测地曲率的符号. 本册的整体 Gauss--Bonnet 对定向曲面陈述; 闭嵌入曲面可定向是明确使用的拓扑事实.
\begin{example}
计算球面与圆柱的度量和面积元, 并写出圆柱展开.
\end{example}
\begin{solution}
球面参数如前, 在 $0<\theta<\pi$ 上 $E=R^2,F=0,G=R^2\sin^2\theta$. 因而 $\dd A=R^2\sin\theta\dd\theta\dd\varphi$, 对 $0<\varphi<2\pi$ 积分得 $4\pi R^2$, 省略的极点和经线面积为0.

圆柱参数 $X(u,v)=(R\cos(u/R),R\sin(u/R),v)$ 有 $E=G=1,F=0$, 故与平面 $(u,v)$ 的度量相同. 限制 $u$ 于长度小于 $2\pi R$ 的区间得到单射展开. 若不限制, $u$ 与 $u+2\pi R$ 对应同一点, 这就是全局展开的重数问题.
\end{solution}
\originalsection{习题}
\begin{exercise}\label{ex:4a}
对图像 $X(u,v)=(u,v,f(u,v))$, 求第一基本形式和面积元. 用它计算抛物面 $z=(x^2+y^2)/2$ 在 $x^2+y^2\le a^2$ 上的面积.
\end{exercise}
\begin{exercise}\label{ex:4b}
对球面去掉北极, 采用参数 $X(u,v)=\bigl(2u,2v,u^2+v^2-1\bigr)/(1+u^2+v^2)$. 证明其为单位球面参数, 求度量并判断它是否保角.
\end{exercise}
\begin{exercise}\label{ex:4c}
圆锥去掉顶点, 参数 $X(r,\theta)=(r\cos\theta,r\sin\theta,ar)$, $a>0,r>0$. 构造局部等距到平面扇形的参数变换, 并给出绕锥一周对应的扇形角.
\end{exercise}

\originalchapter{第二基本形式与曲率}
\originalsection{形算子与第二基本形式}
第一基本形式只测量曲面内的长度. 曲面在空间中如何弯曲, 需要比较邻点的法向量. 本章选定单位法向 $N$; 所有导数首先作为欧氏空间导数计算.
\begin{definition}
形算子 $S_p:T_pM\to T_pM$ 为 $S_p(V)=-\dd N_p(V)$. 第二基本形式为 $\II_p(V,W)=\langle S_pV,W\rangle$.
\end{definition}
因 $|N|^2=1$, $\dd N(V)\perp N$, 故它确实属于 $T_pM$. 在参数坐标中记
\[
 h_{ij}=\langle X_{ij},N\rangle,\qquad
 h=\begin{pmatrix}e&f\\f&g_2\end{pmatrix}.
\]
这里用 $g_2$ 作为第二基本形式的末项, 避免与度量矩阵 $g$ 混淆. 具体算例常写 $\II=e\dd u^2+2f\dd u\dd v+g_2\dd v^2$.
\begin{proposition}\label{prop:shape}
第二基本形式对称, $\langle S X_i,X_j\rangle=h_{ij}$. 形算子的坐标矩阵为 $g^{-1}h$, 并且它关于度量 $g$ 自伴.
\end{proposition}
\begin{proof}
微分 $\langle N,X_j\rangle=0$, 得 $\langle-\partial_iN,X_j\rangle=\langle N,X_{ij}\rangle$. 混合偏导交换说明 $h_{ij}=h_{ji}$. 若 $SX_i=\sum_k S^k{}_iX_k$, 则 $h_{ij}=\sum_k S^k{}_i g_{kj}$. 矩阵形式为 $h=S^{\mathsf T}g$; 利用 $h$ 对称也有 $h=gS$, 故 $S=g^{-1}h$. 对称性正是 $\langle SV,W\rangle=\langle V,SW\rangle$.
\end{proof}
\begin{definition}
形算子的两个实特征值 $k_1,k_2$ 称为主曲率, 正交特征方向称为主方向. Gauss 曲率为 $K=k_1k_2=\det S$, 平均曲率为 $H=(k_1+k_2)/2=\tr S/2$.
\end{definition}
线性代数中的实对称算子谱定理保证主方向存在. $g^{-1}h$ 在一般坐标下未必是普通对称矩阵, 但它与对称矩阵 $g^{-1/2}hg^{-1/2}$ 相似, 因而仍有实谱. 特征方程应写 $\det(h-kg)=0$, 不能写成 $\det(g-kh)=0$.
\begin{proposition}\label{prop:KH}
在参数坐标中,
\[
 K=\frac{eg_2-f^2}{EG-F^2},\qquad
 H=\frac{Eg_2-2Ff+Ge}{2(EG-F^2)}.
\]
反转法向量时 $h,S,k_i,H$ 全部变号, 而 $K$ 不变.
\end{proposition}
\begin{proof}
第一式是 $\det(g^{-1}h)=\det h/\det g$. 用 $g^{-1}=(EG-F^2)^{-1}\left(\begin{smallmatrix}G&-F\\-F&E\end{smallmatrix}\right)$ 求迹得第二式. 反向结论来自 $N\mapsto-N$.
\end{proof}
\originalsection{法曲率与 Euler 公式}
\begin{definition}
单位切向量 $V$ 的法曲率为 $k_n(V)=\II(V,V)$. 对非零切向量可写成 $k_n(V)=\II(V,V)/\I(V,V)$.
\end{definition}
\begin{proposition}\label{prop:normal-curvature}
任意经过 $p$、单位速度且切向为 $V$ 的曲面曲线 $\alpha$, 都满足 $\langle\alpha''(0),N(p)\rangle=k_n(V)$. 若 $V=\cos\theta e_1+\sin\theta e_2$, $e_i$ 为单位主方向, 则
\begin{equation}\label{eq:euler}
 k_n(V)=k_1\cos^2\theta+k_2\sin^2\theta.
\end{equation}
\end{proposition}
\begin{proof}
微分 $\langle\alpha',N\rangle=0$, 得 $\langle\alpha'',N\rangle=-\langle\alpha',\dd N(\alpha')\rangle=\II(V,V)$. 展开自伴算子在正交特征基下的二次型, 混合项为0, 得 Euler 公式.
\end{proof}
所以法曲率只依赖切向, 不依赖曲线如何在曲面上偏转. 若空间曲率 $\kappa\ne0$, 曲线主法向与曲面法向的夹角为 $\beta$, 则 $k_n=\kappa\cos\beta$; 这就是 Meusnier 关系. 法截线的空间曲率为 $|k_n|$; 当 $k_n\ne0$ 时, 其密切平面是 $V,N$ 所张成的平面. 在渐近方向上 $k_n=0$, 该点的 Frenet 密切平面不一定有定义.
\begin{definition}
点的两主曲率相等时称为脐点. $K>0$ 称为椭圆点, $K<0$ 称为双曲点. $K=0$ 且 $S\ne0$ 称为抛物点, $S=0$ 称为平坦点. $\II(V,V)=0$ 的非零方向称为渐近方向.
\end{definition}
Euler 公式说明: 椭圆点无法向曲率为0的方向; 双曲点有两条不同渐近方向; 抛物点只有一条. 平坦点在此是点态概念, 不自动意味着邻域是平面.
\begin{center}
\begin{tikzpicture}
\begin{axis}[width=6.6cm,height=4.4cm,view={40}{25},axis lines=box,clip=false,xtick=\empty,ytick=\empty,ztick=\empty]
\addplot3[surf,domain=-1:1,y domain=-1:1,samples=13,shader=flat,z buffer=sort,colormap={dg}{color(0cm)=(white);color(1cm)=(structurecolor)},opacity=.85] {(.5*x*x-.5*y*y)};
\node[anchor=north] at (axis description cs:.35,-.05) {$u$};
\node[anchor=west] at (axis description cs:.92,.02) {$v$};
\node[anchor=east] at (axis description cs:-.03,.50) {$z$};
\end{axis}
\end{tikzpicture}
\end{center}
在鞍面 $z=(u^2-v^2)/2$ 的原点, $\I=\dd u^2+\dd v^2$, $\II=\dd u^2-\dd v^2$, $K=-1$, 两个渐近方向为 $u=\pm v$. 图按所给鞍面绘制, 不是曲率证明.
\originalsection{基本算例与脐点曲面}
\begin{example}
计算平面、半径 $R$ 的球面和圆柱的主曲率.
\end{example}
\begin{solution}
平面法向量固定, $S=0$. 球面外法向 $N(p)=p/R$, 故对任意切向 $V$, $SV=-V/R$, 两主曲率为 $-1/R$, $K=1/R^2$, $H=-1/R$.

圆柱 $X(\theta,z)=(R\cos\theta,R\sin\theta,z)$ 取外法向 $(\cos\theta,\sin\theta,0)$. 环向单位向量 $e_\theta=(-\sin\theta,\cos\theta,0)$ 满足 $Se_\theta=-e_\theta/R$, 轴向向量满足 $Se_z=0$. 所以 $K=0$, $H=-1/(2R)$. 圆柱和其局部展开平面有相同第一基本形式, 却有不同第二基本形式.
\end{solution}
\begin{proposition}\label{prop:graph-curvature}
图像曲面 $X(u,v)=(u,v,f(u,v))$ 取向上法向, 令 $W=\sqrt{1+f_u^2+f_v^2}$, 则
\[
 \II=\frac{f_{uu}\dd u^2+2f_{uv}\dd u\dd v+f_{vv}\dd v^2}{W},\qquad
 K=\frac{f_{uu}f_{vv}-f_{uv}^2}{(1+f_u^2+f_v^2)^2}.
\]
\end{proposition}
\begin{proof}
$N=(-f_u,-f_v,1)/W$, $X_{ij}=(0,0,f_{ij})$, 得第二基本形式. 第一基本形式行列式为 $W^2$, 而 $\det h=(f_{uu}f_{vv}-f_{uv}^2)/W^2$, 相除得 $K$.
\end{proof}
\begin{theorem}\label{thm:umbilic}
连通正则曲面若处处为脐点, 则它的像包含在某个平面或球面中.
\end{theorem}
\begin{proof}
在一个连通参数片上写 $S=\lambda\id$, 即 $N_u=-\lambda X_u$, $N_v=-\lambda X_v$. 对两式取混合偏导并相减, 得 $\lambda_vX_u=\lambda_uX_v$. 两向量独立, 故 $\lambda_u=\lambda_v=0$, $\lambda$ 局部恒定. 连通性使同一定向下它恒定; 若整体不先选定定向, 在重叠片上法向符号可统一延续, 以下局部平面或球面也因非空开交集而必须一致.

若 $\lambda=0$, $N$ 固定, $\langle X,N\rangle$ 固定, 曲面落在平面上. 若 $\lambda\ne0$, $X+N/\lambda$ 的两偏导为0, 圆心固定且半径为 $1/|\lambda|$. 不同片若给出两不同平面或球面, 其交不能含二维开片, 因而沿连通曲面得到同一个平面或球面.
\end{proof}
\originalsection{习题}
\begin{exercise}\label{ex:5a}
求旋转抛物面 $z=(x^2+y^2)/(2a)$, $a>0$, 的 $K,H$, 取向上法向. 判断其椭圆点和脐点.
\end{exercise}
\begin{exercise}\label{ex:5b}
设一点 $k_1=2,k_2=-1$. 求所有法曲率为0的方向, 并求法曲率的最大值、最小值和方向平均值.
\end{exercise}
\begin{exercise}\label{ex:5c}
正则水平曲面 $F=0$ 取 $N=\nabla F/|\nabla F|$. 证明 $\II(V,W)=-\operatorname{Hess}F(V,W)/|\nabla F|$ 对切向量成立. 据此求椭球 $x^2/a^2+y^2/b^2+z^2/c^2=1$, $a,b,c>0$, 在 $(a,0,0)$ 的主曲率.
\end{exercise}

\originalchapter{旋转曲面与极小曲面}
\originalsection{旋转曲面的基本形式}
旋转曲面将一个平面母线绕轴旋转, 是同时练习基本形式、曲率和测地线的贯穿算例. 令母线 $(r(s),z(s))$ 以弧长参数表示, 满足 $r>0$, $r'^2+z'^2=1$, 参数片为
\[
 X(s,\theta)=(r(s)\cos\theta,r(s)\sin\theta,z(s)).
\]
\begin{proposition}\label{prop:rotation}
取 $N=(-z'\cos\theta,-z'\sin\theta,r')$, 则
\[
 \I=\dd s^2+r^2\dd\theta^2,\qquad
 \II=(-r''z'+z''r')\dd s^2+rz'\dd\theta^2.
\]
主方向为经向与纬向, 主曲率为 $k_s=-r''z'+z''r'$, $k_\theta=z'/r$, 且 $K=-r''/r$.
\end{proposition}
\begin{proof}
$X_s=(r'\cos\theta,r'\sin\theta,z')$, $X_\theta=(-r\sin\theta,r\cos\theta,0)$, 两者正交且长度为 $1,r$. 其叉积归一化即为所给 $N$. 二阶导数取法向分量得到 $e=-r''z'+z''r'$, $f=0$, $g_2=rz'$, 所以两坐标方向为主方向.

微分单位速度条件得 $r'r''+z'z''=0$. 将 $e z'$ 展开为 $-r''z'^2+z''r'z'=-r''(z'^2+r'^2)=-r''$. 因而 $K=k_sk_\theta=-r''/r$. 此推导不除以 $z'$, 在 $z'=0$ 的点同样有效.
\end{proof}
例如球面母线 $r=R\sin(s/R),z=R\cos(s/R)$ 给出 $K=1/R^2$; 圆柱母线 $r=R,z=s$ 给出 $K=0$. 内蕴曲率只由 $r$ 的二阶变化决定, 而各主曲率仍依赖法向选择.
\begin{example}
求标准环面 $X(\theta,\varphi)=((R+a\cos\varphi)\cos\theta,(R+a\cos\varphi)\sin\theta,a\sin\varphi)$ 的面积与 Gauss 曲率, $R>a>0$.
\end{example}
\begin{solution}
直接有
\[
 \I=(R+a\cos\varphi)^2\dd\theta^2+a^2\dd\varphi^2,\qquad
 \dd A=a(R+a\cos\varphi)\dd\theta\dd\varphi.
\]
取指向管外的法向 $N=(\cos\varphi\cos\theta,\cos\varphi\sin\theta,\sin\varphi)$, 两主曲率为 $-1/a$ 和 $-\cos\varphi/(R+a\cos\varphi)$. 因而
\[
 K=\frac{\cos\varphi}{a(R+a\cos\varphi)},\qquad \Area=4\pi^2aR.
\]
环面外侧 $K>0$, 内侧 $K<0$, 上下两条纬圆 $K=0$. 积分 $K\dd A=\cos\varphi\dd\theta\dd\varphi$ 得总 Gauss 曲率0, 后面的 Gauss--Bonnet 将解释该值与形状细节无关.
\end{solution}
\begin{center}
\begin{tikzpicture}[scale=2]
\begin{axis}[width=8.6cm,height=5.2cm,view={35}{30},axis lines=none,axis equal image]
\addplot3[surf,domain=0:360,y domain=0:360,samples=25,samples y=17,shader=flat,z buffer=sort,colormap={ring}{color(0cm)=(white);color(1cm)=(structurecolor)}] ({(2+.65*cos(y))*cos(x)},{(2+.65*cos(y))*sin(x)},{.65*sin(y)});
\end{axis}
\end{tikzpicture}
\end{center}
\originalsection{直纹曲面与可展性}
\begin{definition}
由直线族组成的参数片 $X(s,t)=c(s)+td(s)$, $|d(s)|=1$, 称为直纹曲面. 正则部分若 $K=0$, 称为局部可展的直纹曲面; 这里先以曲率条件使用术语, 其局部平直性会在内蕴章节说明.
\end{definition}
\begin{proposition}\label{prop:ruled}
直纹参数片在正则点有
\[
 K=-\frac{\det(c',d,d')^2}{|(c'+td')\times d|^4}\le0.
\]
特别 $K=0$ 当且仅当 $c',d,d'$ 线性相关.
\end{proposition}
\begin{proof}
$X_s=c'+td'$, $X_t=d$, $X_{tt}=0$. 因而第二基本形式 $g_2=0$, $f=\langle d',N\rangle$. 记 $W=|X_s\times d|$, 则
\[
 f=\frac{\langle d',(c'+td')\times d\rangle}{W}
   =\frac{\det(c',d,d')}{W}.
\]
由 $EG-F^2=W^2$, 得 $K=-f^2/W^2$.
\end{proof}
平面、圆柱、圆锥的正则部分均满足零曲率. 另一例子是空间曲线的切线曲面.
\begin{example}
单位速度曲线 $c$ 有 $\kappa>0$, 考察 $X(s,t)=c(s)+tT(s)$ 在 $t>0$ 的正则部分, 求其曲率.
\end{example}
\begin{solution}
$X_s=T+t\kappa n$, $X_t=T$, 叉积为 $-t\kappa b$, 所以 $N=-b$, $E=1+t^2\kappa^2,F=1,G=1$. 有 $X_{st}=\kappa n$, 其法向分量为0, $X_{tt}=0$, 而 $X_{ss}$ 的 $b$ 分量为 $t\kappa\tau$, 故 $e=-t\kappa\tau$. 于是 $K=0$, $H=-\tau/(2t\kappa)$. $t=0$ 时叉积为0, 该参数片在回归边上奇异, 不能把公式延伸至那里.
\end{solution}
\originalsection{极小曲面的面积变分}
\begin{definition}
平均曲率 $H\equiv0$ 的曲面称为极小曲面. 名称指面积一阶变分为0, 不自动保证每一个有限区域都是全局面积最小.
\end{definition}
\begin{theorem}\label{thm:first-area}
取参数区域 $D$ 满足 $\overline D$ 紧含于光滑参数片的定义域, 并取 $f\in C_c^\infty(D)$, 考察法向变分 $X_t=X+t fN$. 当 $t$ 足够小时它保持正则, 且
\[
 \left.\frac{\dd}{\dd t}\Area(X_t(D))\right|_{t=0}
 =-2\int_D fH\dd A.
\]
在此面积按参数片积分计, 对单射片等于实际面积. 对所有这种变分一阶导数为0, 当且仅当 $H=0$.
\end{theorem}
\begin{proof}
一阶导数 $\partial_iX_t=X_i+t(f_iN+fN_i)$. 故度量变分为 $\dot g_{ij}=f(\langle N_i,X_j\rangle+\langle X_i,N_j\rangle)=-2fh_{ij}$. 对正定矩阵有 $\partial_t\det g=(\det g)\tr(g^{-1}\dot g)$, 这可由行列式的多线性或伴随矩阵公式得到. 因而
\[
 \partial_t\sqrt{\det g}\big|_0
 =\tfrac12\sqrt{\det g}\tr(-2fg^{-1}h)=-2fH\sqrt{\det g}.
\]
紧支撑使积分微分合法且边界固定. 若积分对任意 $f$ 为0, 在 $H$ 非零点附近取非负、非零紧支撑函数, 使 $H$ 在其支撑上保持同号会得到非零积分, 矛盾. 正则性在紧支撑集上由原叉积的正下界和连续性保持, 支撑外曲面不变.
\end{proof}
\begin{proposition}
图像 $z=f(u,v)$ 为极小曲面当且仅当
\begin{equation}\label{eq:minimal-graph}
 (1+f_v^2)f_{uu}-2f_uf_vf_{uv}+(1+f_u^2)f_{vv}=0.
\end{equation}
\end{proposition}
\begin{proof}
将图像的第一、第二基本形式代入命题\ref{prop:KH}, 分母为 $2(1+f_u^2+f_v^2)^{3/2}>0$, 分子恰为左侧, 所以等价. 这也等价于 $\operatorname{div}(\nabla f/\sqrt{1+|\nabla f|^2})=0$.
\end{proof}
\begin{example}
验证悬链面和螺旋面都是极小曲面, 并计算 Gauss 曲率.
\end{example}
\begin{solution}
悬链面取 $X(u,v)=(a\cosh v\cos u,a\cosh v\sin u,av)$, $a>0$. 有 $E=G=a^2\cosh^2v,F=0$. 选 $N=(\sech v\cos u,\sech v\sin u,-\tanh v)$, 二阶法向系数为 $e=-a,f=0,g_2=a$. 两主曲率为 $\mp1/(a\cosh^2v)$, 故 $H=0$, $K=-1/(a^2\cosh^4v)$.

螺旋面 $Y(u,v)=(u\cos v,u\sin v,av)$ 的 $E=1,F=0,G=u^2+a^2$. 取 $N=(a\sin v,-a\cos v,u)/\sqrt{u^2+a^2}$. 则 $e=g_2=0,f=-a/\sqrt{u^2+a^2}$, 故 $H=0$, $K=-a^2/(u^2+a^2)^2$. 两曲面都具有负 Gauss 曲率, 所以极小不等于平直.
\end{solution}
\originalsection{习题}
\begin{exercise}\label{ex:6a}
环面上取环向与管向单位主方向. 当 $R=2a$, 求环面上满足 $H=0$ 的点集, 并说明整个环面是否极小.
\end{exercise}
\begin{exercise}\label{ex:6b}
圆锥 $X(r,\theta)=(r\cos\theta,r\sin\theta,ar)$, $a>0$, $r>0$. 求两主曲率和平均曲率, 取 $X_r\times X_\theta$ 确定的法向.
\end{exercise}
\begin{exercise}\label{ex:6c}
证明没有非空紧致无边界极小嵌入曲面 $M\subset\R^3$. 提示可考察离固定原点距离最大的点, 但须写出二阶导数的符号约束.
\end{exercise}

\par\Needspace{9\baselineskip}
\originalchapter{内蕴联络与 Gauss 方程}
\originalsection{切向投影与 Christoffel 符号}
欧氏导数会离开曲面切平面. 若只保留切向分量, 得到曲面自身的微分法则. 设 $V,W$ 为切向量场, 用 $D_VW$ 表示沿 $V$ 对三维向量分量求导.
\begin{definition}
曲面的协变导数为 $\nabla_VW=(D_VW)^{\top}$, 其中 $\top$ 表示正交投影到切平面. 法向分解为
\begin{equation}\label{eq:gauss-formula}
 D_VW=\nabla_VW+\II(V,W)N.
\end{equation}
\end{definition}
式\eqref{eq:gauss-formula}的法向系数来自微分 $\langle W,N\rangle=0$. 协变导数对第一个变量是函数线性的, 对第二个变量满足
\[
 \nabla_V(fW)=V(f)W+f\nabla_VW.
\]
它同时满足度量相容性
\[
 V\langle W,Z\rangle=\langle\nabla_VW,Z\rangle+\langle W,\nabla_VZ\rangle
\]
及无挠性 $\nabla_VW-\nabla_WV=[V,W]$, 其中 Lie 括号在坐标中是向量场作为微分算子的交换子. 两性质分别继承自欧氏内积求导和混合偏导交换.
\begin{definition}
坐标切向量记为 $X_i$, 定义 Christoffel 符号 $\Gamma^k_{ij}$ 使
\[
 X_{ij}=\sum_k\Gamma^k_{ij}X_k+h_{ij}N.
\]
\end{definition}
\begin{proposition}\label{prop:christoffel}
记 $(g^{ij})=g^{-1}$, 则
\begin{equation}\label{eq:christoffel}
 \Gamma^k_{ij}=\frac12\sum_\ell g^{k\ell}
 (\partial_i g_{j\ell}+\partial_j g_{i\ell}-\partial_\ell g_{ij}).
\end{equation}
所以联络由第一基本形式决定.
\end{proposition}
\begin{proof}
微分 $g_{j\ell}=\langle X_j,X_\ell\rangle$ 并写出指标循环的三个式子, 得
\[
 \partial_i g_{j\ell}+\partial_j g_{i\ell}-\partial_\ell g_{ij}
 =2\langle X_{ij},X_\ell\rangle
 =2\sum_k\Gamma^k_{ij}g_{k\ell}.
\]
左乘逆矩阵即得. 对 $W=\sum_jw^jX_j$, 分量公式为
\[
 \nabla_{X_i}W=\sum_k\left(\partial_iw^k+\sum_j\Gamma^k_{ij}w^j\right)X_k,
\]
故知道 $g$ 便知道整个联络, 而无须知道嵌入的第二基本形式.
\end{proof}
Christoffel 符号在换坐标时含二阶换元项, 所以它们本身不是一个张量的分量. 由它们定义的协变导数却是坐标无关的几何操作. 用同一公式, 对不一定嵌入三维空间的二维正定光滑度量也能定义联络; 后面的双曲平面使用这一做法.
\originalsection{Gauss 曲率的内蕴性}
为避免符号混乱, 固定曲率算子的约定
\[
 R(V,W)Z=\nabla_V\nabla_WZ-\nabla_W\nabla_VZ-\nabla_{[V,W]}Z.
\]
\begin{theorem}[Gauss 方程]\label{thm:gauss}
对曲面切向量有
\begin{equation}\label{eq:gauss-eq}
 R(V,W)Z=\II(W,Z)SV-\II(V,Z)SW.
\end{equation}
对任意正向单位正交切标架 $e_1,e_2$,
\[
 K=\langle R(e_1,e_2)e_2,e_1\rangle.
\]
\end{theorem}
\begin{proof}
欧氏空间的导数交换满足 $D_VD_WZ-D_WD_VZ-D_{[V,W]}Z=0$. 将两次导数都按式\eqref{eq:gauss-formula}分解. 其中 $D_V(\II(W,Z)N)$ 的切向分量为 $-\II(W,Z)SV$. 收集切向分量得
\[
 0=R(V,W)Z-\II(W,Z)SV+\II(V,Z)SW,
\]
即所需公式. 再令 $V=e_1,W=Z=e_2$, 内积为 $h_{22}h_{11}-h_{12}^2=\det S=K$. 并不要求此标架是主标架.
\end{proof}
\begin{corollary}[Gauss 绝妙定理]\label{cor:egregium}
Gauss 曲率只依赖第一基本形式及其一、二阶导数. 局部等距保持 Gauss 曲率.
\end{corollary}
\begin{proof}
式\eqref{eq:christoffel}表明协变导数由度量决定; 曲率算子由协变导数、它的一阶导数和 Lie 括号决定. Gauss 方程把 $K$ 表示成这个内蕴算子的一个度量分量. 局部等距在相应坐标中保持度量, 因而保持该值.
\end{proof}
这解释了为何圆柱可局部展开而球面不能: 平面 $K=0$, 球面 $K=1/R^2$. 但相同 $K$ 未必给出相同局部度量, 更不能由一个数恢复嵌入曲面.
\begin{proposition}[Codazzi 方程]\label{prop:codazzi}
形算子满足 $(\nabla_VS)W=(\nabla_WS)V$, 其中 $(\nabla_VS)W=\nabla_V(SW)-S(\nabla_VW)$. 等价地, 第二基本形式的协变导数对前两指标对称.
\end{proposition}
\begin{proof}
对法向场使用欧氏导数交换式. 因 $D_WN=-SW$, 切向分量给出 $-\nabla_V(SW)+\nabla_W(SV)+S[V,W]=0$. 用无挠性替换 $[V,W]=\nabla_VW-\nabla_WV$, 即得结论. 对任意 $Z$ 取内积并用度量相容性, 得第二基本形式版本.
\end{proof}
\originalsection{正交标架与曲率二形式}
二维上用一条连接一形式能把许多指标浓缩. 在局部选正向单位正交标架 $e_1,e_2$, 定义一形式 $\omega$ 使
\[
 \nabla_Ve_1=\omega(V)e_2,\qquad \nabla_Ve_2=-\omega(V)e_1.
\]
这里 $\dd A$ 作为二形式表示与曲面定向相容的面积形式. 一形式是对切向量线性取值的场, 可在坐标中写 $\omega=A\dd u+B\dd v$. 其外微分在本册只需二元公式 $\dd\omega=(\partial_uB-\partial_vA)\dd u\wedge\dd v$; Green 公式为 $\int_D\dd\omega=\int_{\partial D}\omega$.
\begin{proposition}\label{prop:connection-curvature}
在上述约定下, $\dd\omega=-K\dd A$. 若把标架旋转为 $\widetilde e_1=\cos\psi e_1+\sin\psi e_2$, $\widetilde e_2=-\sin\psi e_1+\cos\psi e_2$, 则 $\widetilde\omega=\omega+\dd\psi$.
\end{proposition}
\begin{proof}
计算 $R(V,W)e_1$ 的 $e_2$ 分量, 两个连接乘积相消, 只剩 $V(\omega(W))-W(\omega(V))-\omega([V,W])=\dd\omega(V,W)$. 对 $V=e_1,W=e_2$, Gauss 方程给出 $R(e_1,e_2)e_1$ 的 $e_2$ 分量为 $-K$. 而 $\dd A(e_1,e_2)=1$, 得首式. 对旋转后的两向量直接用乘积法则微分, 新增角速度项 $\dd\psi$, 得第二式.
\end{proof}
\begin{proposition}\label{prop:metric-K}
若 $\I=\dd u^2+a(u,v)^2\dd v^2$, $a>0$, 则 $K=-a_{uu}/a$. 若 $\I=e^{2f}(\dd u^2+\dd v^2)$, 则 $K=-e^{-2f}(f_{uu}+f_{vv})$.
\end{proposition}
\begin{proof}
第一种度量取 $e_1=\partial_u$, $e_2=a^{-1}\partial_v$. 由 Christoffel 公式, $\omega(\partial_u)=0$, $\omega(\partial_v)=a_u$, 所以 $\dd\omega=a_{uu}\dd u\wedge\dd v$. 面积元为 $a\dd u\dd v$, 得公式.

第二种取 $e_1=e^{-f}\partial_u$, $e_2=e^{-f}\partial_v$. Christoffel 公式给出 $\omega=-f_v\dd u+f_u\dd v$, 从而 $\dd\omega=(f_{uu}+f_{vv})\dd u\wedge\dd v$. 面积元为 $e^{2f}\dd u\dd v$, 亦得公式. 两次计算使用同一符号约定.
\end{proof}
\originalsection{曲面的基本数据}
\begin{theorem}[刚性部分]\label{thm:surface-rigidity}
若两个正则参数片 $X,\widetilde X$ 在同一连通参数域上具有相同的第一、第二基本形式, 法向均由相容参数方向选取, 则它们相差一个正向欧氏刚体运动.
\end{theorem}
\begin{proof}
在每点以 $(X_1,X_2,N)$ 和 $(\widetilde X_1,\widetilde X_2,\widetilde N)$ 组成可逆矩阵 $F,\widetilde F$. Gauss 公式与 $N_i=-\sum_k(g^{-1}h)^k{}_iX_k$ 写成 $\partial_iF=F B_i$, 其中 $B_i$ 只由 $g,h$ 与 Christoffel 符号组成. 两个片的 $B_i$ 相同, 所以 $A=\widetilde FF^{-1}$ 的两偏导为0, 在连通域上恒定. 因 $F^{\mathsf T}F=\widetilde F^{\mathsf T}\widetilde F=\operatorname{diag}(g,1)$, $A$ 为正交矩阵. 两组框架正向, 故 $\det A=1$. 又 $\partial_i(\widetilde X-AX)=0$, 平移向量恒定.
\end{proof}
\begin{theorem}[曲面基本定理的局部存在性]\label{thm:surface-existence}
在平面开集上给定光滑正定对称矩阵 $g$ 与光滑对称矩阵 $h$. 用 $g$ 定义联络, 用 $S=g^{-1}h$ 定义形算子. 若它们满足式\eqref{eq:gauss-eq}与 Codazzi 方程, 则在每点的足够小邻域上存在正则曲面参数片, 第一、第二基本形式恰为 $g,h$. 相同参数域上的实现由前一定理在刚体运动意义下唯一.
\end{theorem}
\begin{proof}
取以该点为原点的小矩形, 指标 $i,j,k$ 取1,2. 记 $\Gamma_i$ 的第 $k,j$ 项为 $\Gamma^k_{ij}$, $S_i$ 为 $S$ 的第 $i$ 列, $h_i$ 为 $h$ 的第 $i$ 行. 设置三阶矩阵
\[
 B_i=\begin{pmatrix}\Gamma_i&-S_i\\h_i&0\end{pmatrix},\qquad
 M=\begin{pmatrix}g&0\\0&1\end{pmatrix}.
\]
度量相容性及 $gS=h$ 给出 $\partial_iM=B_i^{\mathsf T}M+MB_i$. Gauss--Codazzi 使
\begin{equation}\label{eq:flat-system}
 \partial_1B_2-\partial_2B_1+B_1B_2-B_2B_1=0.
\end{equation}
具体而言, 左上块为联络曲率矩阵减去 $S_1h_2-S_2h_1$, 正是 Gauss 方程; 左下块为 $\partial_1h_2-\partial_2h_1+h_1\Gamma_2-h_2\Gamma_1$, 是 Codazzi 的坐标式. 右上块是该式在度量相容性下的转置对偶, 右下块为 $-h_1S_2+h_2S_1=0$, 来自 $hg^{-1}h$ 对称. 因而所有块确实为0.

选择初始可逆框架矩阵 $F_0$ 满足 $F_0^{\mathsf T}F_0=M(0)$ 且 $\det F_0>0$. 先沿 $v=0$ 解 $F_u=FB_1$, 再从每个底边点沿竖线解 $F_v=FB_2$. 光滑常微分方程解依赖参数, 所得 $F$ 光滑. 记误差 $H=F_u-FB_1$. 对 $v$ 求导并用式\eqref{eq:flat-system}, 得 $H_v=HB_2$, 底边 $H=0$, 故 $H$ 恒为0. 这样两个偏微分方程都满足, 不必未经说明引用路径无关积分定理.

矩阵 $F^{\mathsf T}F$ 与 $M$ 在每条上述积分线都满足相同线性矩阵方程和初值, 因而相等. 把 $F$ 的列记为 $Y_1,Y_2,N$, 则 $\langle Y_i,Y_j\rangle=g_{ij}$, $N$ 为单位法向. 正定性使 $Y_1,Y_2$ 独立, 连续性保持正向. 方程的前两列写成 $\partial_iY_j=\sum_k\Gamma^k_{ij}Y_k+h_{ij}N$. $\Gamma^k_{ij}$ 与 $h_{ij}$ 均对 $i,j$ 对称, 所以 $\partial_1Y_2=\partial_2Y_1$.

定义 $X(u,v)=p_0+\int_0^uY_1(t,0)\dd t+\int_0^vY_2(u,t)\dd t$. 有 $X_v=Y_2$, 而 $X_u=Y_1(u,0)+\int_0^v\partial_uY_2(u,t)\dd t=Y_1(u,v)$. 因而 $X$ 正则, 度量为 $g$, 二阶法向分量为 $h$, 所选 $N$ 恰与参数定向相容. 这完成局部存在性.
\end{proof}
局部积分不保证任意全局数据都能拼成一个单射嵌入; 非单连通区域可能出现框架或位置的回路不一致. 本定理精确到足够小邻域, 不把局部实现扩张成未经证明的全局嵌入结论.
\originalsection{习题}
\begin{exercise}\label{ex:7a}
对 $\I=\dd u^2+r(u)^2\dd v^2$, 写出全部非零 Christoffel 符号, 并由公式直接求 $K$.
\end{exercise}
\begin{exercise}\label{ex:7b}
给定 $\I=e^{2f}(\dd u^2+\dd v^2)$ 且 $f(u,v)=u^2+v^2$, 求 $K$ 和连接形式. 此度量是否局部等距于平面?
\end{exercise}
\begin{exercise}\label{ex:7c}
假设一个连通参数片 $K\equiv0$. 证明在足够小的单连通邻域可以选择平行正交标架, 并构造到欧氏平面的局部等距. 须说明为何所用一形式可积分.
\end{exercise}

\originalchapter{测地线与变分}
\originalsection{测地线方程}
曲面上“直行”不能用欧氏加速度为0定义, 因为留在曲面上通常需要法向约束力. 应要求加速度没有切向分量.
\begin{definition}
光滑曲线 $\gamma:I\to M$ 若 $\nabla_{\dot\gamma}\dot\gamma=0$, 称为测地线. 对嵌入曲面, 等价于 $\ddot\gamma$ 垂直于切平面. 允许常值测地线.
\end{definition}
\begin{proposition}\label{prop:geodesic-eq}
在坐标 $x=(x^1,x^2)$ 中, 测地线方程为
\begin{equation}\label{eq:geodesic}
 \ddot x^k+\sum_{i,j}\Gamma^k_{ij}(x)\dot x^i\dot x^j=0,\qquad k=1,2.
\end{equation}
测地线速度恒定. 对任意 $p\in M,V\in T_pM$, 存在包含0的开区间及唯一测地线满足 $\gamma(0)=p,\dot\gamma(0)=V$; 唯一性指共同定义区间上相同. 解光滑依赖初值.
\end{proposition}
\begin{proof}
把 $\dot\gamma=\sum_i\dot x^iX_i$ 代入协变导数分量公式, 得方程. 度量相容性给出 $(|\dot\gamma|^2)'=2\langle\nabla_{\dot\gamma}\dot\gamma,\dot\gamma\rangle=0$.

将二阶方程写成一阶系统 $\dot x=v$, $\dot v^k=-\sum\Gamma^k_{ij}(x)v^iv^j$. 光滑右端在坐标域内局部 Lipschitz, 标准常微分方程定理给出存在唯一性及初值的光滑依赖. 换坐标后方程由同一个几何条件描述, 故局部解在交集上相容, 可延续到最大区间. 此处不把局部存在误写成所有曲面上的全时间存在.
\end{proof}
非恒定测地线的仿射换参 $t\mapsto at+b$ 仍是测地线; 一般非匀速换参则增添沿切向的加速度, 不再满足该方程. 一个曲线的像可为测地线轨迹, 但所给参数可能不是测地线参数.
\begin{example}
求半径 $R$ 球面与圆柱上的全部测地线.
\end{example}
\begin{solution}
球面上单位速度测地线满足 $\gamma''=\lambda\gamma$. 对 $|\gamma|^2=R^2$ 微分两次得 $\langle\gamma'',\gamma\rangle=-1$, 因而 $\lambda=-1/R^2$. 初值 $p,V$, $|V|=1$, $p\perp V$ 的解是 $\gamma(s)=p\cos(s/R)+RV\sin(s/R)$, 即大圆. 任意恒速情形用仿射换参, 再加常值曲线即为全部.

圆柱展开坐标 $(u,z)$ 下度量为 $\dd u^2+\dd z^2$, Christoffel 符号全为0, 解为 $u=at+b,z=ct+d$. 在空间中即 $\gamma(t)=(R\cos((at+b)/R),R\sin((at+b)/R),ct+d)$, 包括母线、水平圆及圆柱螺线. 这些曲线都局部直行, 水平圆走过半圈后却不再是两端点间最短的弧.
\end{solution}
\originalsection{能量与第一变分}
\begin{definition}
固定参数区间 $[a,b]$, 曲线能量为 $\mathcal E(\gamma)=\frac12\int_a^b|\dot\gamma|^2\dd t$. 光滑变分为映射 $F(t,\eps)$, 每条 $F(\cdot,\eps)$ 留在曲面上, $F(t,0)=\gamma(t)$. 变分场为 $V(t)=\partial_\eps F(t,0)$.
\end{definition}
\begin{theorem}[能量第一变分]\label{thm:first-energy}
有
\begin{equation}\label{eq:first-energy}
 \left.\partial_\eps\mathcal E(F(\cdot,\eps))\right|_0
 =\langle V,\dot\gamma\rangle\big|_a^b-
 \int_a^b\langle V,\nabla_{\dot\gamma}\dot\gamma\rangle\dd t.
\end{equation}
所以固定端点变分下的驻值曲线恰为测地线.
\end{theorem}
\begin{proof}
积分微分后得到 $\int\langle\nabla_\eps F_t,F_t\rangle\dd t$. 无挠性及参数偏导交换给出 $\nabla_\eps F_t=\nabla_tF_\eps$. 度量相容性允许分部积分, 得所示边界项与积分项. 若端点固定, $V(a)=V(b)=0$, 测地线使变分为0.

反之若对所有固定端点变分一阶导数为0, 在任意短坐标段上可以实现任意紧支撑切向场: 在参数域中取 $x_\eps=x+\eps v$, 再用参数片映入曲面. 若加速度切向分量在某点不为0, 取其附近与该分量同向的小支撑场, 使积分非零, 矛盾. 故切向加速度处处为0.
\end{proof}
\begin{proposition}\label{prop:length-energy}
固定区间上 $L^2\le2(b-a)\mathcal E$, 等号当且仅当速度几乎处处恒定; 光滑时即处处恒定. 非常值正则曲线能量绝对最小当且仅当它长度绝对最小且采用恒速参数.
\end{proposition}
\begin{proof}
对速度函数用 Cauchy--Schwarz, 得 $\bigl(\int v\bigr)^2\le(b-a)\int v^2$. 取等条件是 $v$ 与常数函数成比例. 任意正则曲线可弧长化再线性调整到同一参数区间, 所以相同像和方向的最小可能能量为 $L^2/(2(b-a))$. 若长度不是最小, 取更短曲线并匀速化会得到更小能量; 若速度不恒定, 同一路径匀速化也会降低能量. 反向用不等式直接得到. 允许分段光滑竞争曲线时可逐段弧长化, 连接点不影响长度和能量积分.
\end{proof}
驻值与最小必须区分. 球面上完整大圆是测地线, 但连接同一点的非恒定一周曲线能量大于常值曲线. 前述变分证明只给必要驻值条件; 局部最短性下一章另证.
\originalsection{旋转曲面与 Clairaut 关系}
对第6章的度量 $\I=\dd s^2+r(s)^2\dd\theta^2$, 方程为
\[
 \ddot s-rr'\dot\theta^2=0,\qquad
 \ddot\theta+2(r'/r)\dot s\dot\theta=0.
\]
\begin{proposition}\label{prop:clairaut}
旋转曲面测地线有守恒量 $c=r(s)^2\dot\theta$. 若采用单位速度参数, 与经向的有向夹角为 $\beta$, 则 $c=r(s)\sin\beta$; 同时 $\dot s^2+c^2/r(s)^2=1$.
\end{proposition}
\begin{proof}
第二个测地线方程乘 $r^2$ 后恰为 $(r^2\dot\theta)'=0$. 在单位正交基 $e_s=X_s,e_\theta=X_\theta/r$ 中, 单位切向量为 $\dot s\,e_s+r\dot\theta\,e_\theta$. 其环向分量是 $\sin\beta$, 因而得到 Clairaut 关系, 长度为1给出最后一式.
\end{proof}
若 $c\ne0$, 曲线只能进入 $r\ge|c|$ 的部分; 在 $r=|c|$ 可有 $\dot s=0$ 的转向点. 经线 $\theta$ 恒定始终是测地线. 纬线 $s=s_0$ 在非零速度下为测地线当且仅当 $r'(s_0)=0$.
\begin{example}
判断标准环面的内赤道、外赤道和经过管截面的经圆是否为测地线.
\end{example}
\begin{solution}
用母线弧长 $s=a\varphi$, $r=R+a\cos(s/a)$, 得 $r'=-\sin\varphi$. 在 $\varphi=0,\pi$, $r'=0$, 因而外赤道和内赤道均是恒速测地线. 固定 $\theta$ 的经圆属于经线, 也为测地线. 其他纬线不满足 $r'=0$, 即使它们是空间中的圆, 其加速度仍有曲面切向分量.
\end{solution}
\originalsection{习题}
\begin{exercise}\label{ex:8a}
圆柱半径为 $R$, 两点在展开坐标中为 $(0,0)$ 与 $(R\theta,h)$, $0\le\theta<2\pi$. 列出所有连接两点的测地线, 并确定最短长度和多条最短线出现的条件.
\end{exercise}
\begin{exercise}\label{ex:8b}
旋转曲面单位速度测地线的 Clairaut 常数为 $c\ne0$. 在 $\dot s\ne0$ 的区间求 $\dd\theta/\dd s$, 并说明何时可以在 $r=|c|$ 处把纬线当作测地线.
\end{exercise}
\begin{exercise}\label{ex:8c}
证明若曲面和其度量在一个反射下保持不变, 且该反射的固定点集在曲面上是一条正则曲线, 则该固定曲线采用弧长参数后为测地线. 用初值唯一性, 并说明正则固定曲线条件的作用.
\end{exercise}

\originalchapter{指数映射与最短路径}
\originalsection{指数映射与正常邻域}
初值存在唯一性让我们从一个点沿任意切向出发. 把方向和行进长度合成一个切向量, 得到指数映射.
\begin{definition}
令 $\gamma_V$ 为初值 $p,V$ 的测地线. 对能定义到时刻1的 $V$, 定义 $\Exp_p(V)=\gamma_V(1)$. 由仿射换参, $\Exp_p(tV)=\gamma_V(t)$ 在双方有定义时成立.
\end{definition}
\begin{proposition}\label{prop:normal-neighborhood}
$\Exp_p$ 在切空间原点附近光滑, $\dd(\Exp_p)_0=\id:T_pM\to T_pM$. 因而存在 $\rho>0$, 使它把 $B_\rho(0)\subset T_pM$ 微分同胚地映到一个曲面邻域, 称为正常邻域.
\end{proposition}
\begin{proof}
取足够小的初速度, 常微分方程的局部存在及光滑依赖保证解在 $[0,1]$ 上存在并光滑依赖速度. 对任意 $V$, $\Exp_p(tV)=\gamma_V(t)=p+tV+o(t)$, 因而原点微分是恒等映射. 逆函数定理给出一个开邻域上的微分同胚, 再缩小为切空间球即可.
\end{proof}
不同于一般坐标, 正常坐标把经过中心的径向直线变为测地线. 正常极坐标 $F(r,\theta)=\Exp_p(r e(\theta))$, 其中 $e(\theta)$ 是切平面的单位圆方向, 在 $r=0$ 有极坐标自身的奇异性, 而非曲面奇点.
\begin{theorem}[Gauss 引理]\label{thm:gauss-lemma}
正常极坐标下 $|F_r|=1$, $\langle F_r,F_\theta\rangle=0$. 因而度量写成 $\dd r^2+A(r,\theta)^2\dd\theta^2$, $r>0$, $A>0$.
\end{theorem}
\begin{proof}
每条 $r\mapsto F(r,\theta)$ 是初速度为单位向量的测地线, 故 $|F_r|=1$, $\nabla_rF_r=0$. 微分交叉内积得
\[
 \partial_r\langle F_r,F_\theta\rangle
 =\langle F_r,\nabla_rF_\theta\rangle
 =\langle F_r,\nabla_\theta F_r\rangle
 =\tfrac12\partial_\theta|F_r|^2=0.
\]
$r=0$ 时 $F_\theta=0$, 所以内积恒为0. $r>0$ 时正常坐标的秩为2, $F_\theta$ 非零, 得 $A=|F_\theta|>0$.
\end{proof}
\originalsection{径向测地线的极小性}
\begin{theorem}\label{thm:local-minimizing}
在上述正常邻域, 对 $0<r_0<\rho$, 径向测地线从 $p$ 到 $q=F(r_0,\theta_0)$ 的长度 $r_0$ 在所有曲面分段光滑连接路径中最小. 在邻域内达到等长的正则路径具有同一径向轨迹, 且不回走.
\end{theorem}
\begin{proof}
对留在正常邻域且不经过中心的路径, Gauss 引理给出
\[
 |\dot\gamma|=\sqrt{\dot r^2+A^2\dot\theta^2}\ge|\dot r|.
\]
从中心出发时先忽略 $r<\delta$ 的小段, 对其余段积分, 再令 $\delta\downarrow0$, 得路径长度至少 $r_0$. 多次穿过中心的路径亦可按穿越段分别估计. 径向路径达到此下界.

若竞争路径离开正常邻域, 则在第一次达到任意半径 $r_1\in(r_0,\rho)$ 前已花费长度至少 $r_1>r_0$, 不可能更短. 这是针对整条曲面的竞争路径, 不是只在一个参数域内的比较. 取等要求 $\dot\theta=0$ 在 $r>0$ 处成立, 且径向坐标单调不减; 因 $A>0$, 不允许非零环向速度. 由连续性与正则性, 路径具有所述径向轨迹.
\end{proof}
这只保证足够短的测地线段最短. 指数映射远处可能失去单射性或秩. 球面上从一点到其对跖点有无穷多条长度 $\pi R$ 的大圆半弧, 超过这个长度的径向大圆弧则有更短竞争路径.
\begin{definition}
连通曲面的内蕴距离定义为 $d(p,q)=\inf L(\gamma)$, 下确界取遍从 $p$ 到 $q$ 的分段光滑路径.
\end{definition}
\begin{proposition}\label{prop:distance}
$d$ 是有限值距离, 且诱导曲面原有拓扑. 对嵌入曲面还有 $d(p,q)\ge|p-q|$.
\end{proposition}
\begin{proof}
连通曲面局部路径连通, 路径连通分支既开又闭, 所以分段坐标路径能连接任意两点. 非负性、对称性和三角不等式来自长度、路径反向和拼接. 欧氏空间中积分速度的范数不小于位移范数, 故 $d\ge|p-q|$, 从而不同点距离为正.

对拓扑, 取一个坐标片中包含点的较小闭坐标球, 度量特征值在其中有正下界 $m$ 与上界 $M$. 坐标短直线像的长度不超过 $\sqrt M$ 乘坐标位移, 所以坐标接近推出距离接近. 反之从中心出发要离开此片内小球, 在第一次碰边界前至少花费 $\sqrt m$ 乘球半径的长度. 故充分小的距离球留在任意给定坐标邻域. 两种拓扑一致. 正常邻域定理还给出 $d(p,\Exp_pV)=|V|$ 在充分小的球内成立.
\end{proof}
\originalsection{闭曲面与全局存在}
本节“闭子集”意为在 $\R^3$ 中拓扑闭合, 不等同于紧致无边界. 无边界嵌入曲面可以非紧致而仍为闭子集, 例如平面和完整圆柱.
\begin{theorem}\label{thm:closed-geodesic}
若连通无边界正则嵌入曲面 $M$ 是 $\R^3$ 的闭子集, 则每条测地线可延续到所有实时间, 内蕴距离完备, 任意两点间有恒速长度最小测地线.
\end{theorem}
\begin{proof}
先证完备. $d$-Cauchy 点列由 $|p-q|\le d(p,q)$ 在欧氏空间中也是 Cauchy, 因 $M$ 闭而收敛到 $p\in M$. 由前述拓扑相同, 它也在 $d$ 中收敛. 注意这里不仅用了闭性, 还用了曲面度量与坐标拓扑相容.

对全时间延续, 速度恒定为 $c$. 若最大时间端点 $T<\infty$, 则在任何固定初始时刻之后, 位置留在某个有界欧氏球内. 与 $M$ 相交得到紧集. 速度属于该紧集上长度为 $c$ 的切向量集合, 后者也是紧集: 它是位置与有界速度的闭子集, 切向条件随点连续. 光滑测地线向量场在此紧集附近有有限个局部存在图及统一短时存在界. 于是从接近 $T$ 的时刻再解可越过 $T$, 与最大区间矛盾. 负向时间同理.

最后证最短线存在. 给两点 $p,q$, 取一条连接路径的长度 $L_0$, 取长度趋于 $d(p,q)$ 且不超过 $L_0+1$ 的路径列. 将各路径以恒速参数写在 $[0,1]$ 上. 它们在内蕴距离中 Lipschitz 常数不超过 $L_0+1$, 位置均在欧氏闭球 $\overline B(p,L_0+1)$ 与 $M$ 的紧交集中. 由于两种拓扑一致, 此集在距离 $d$ 下亦紧. 对有理时刻逐次抽取收敛子列, 再取对角子列. 所有曲线有共同 Lipschitz 界, 取有限细时间网格可将网格点收敛提升为整个区间一致收敛. 因此得到连续极限 $\gamma_j\to\gamma$, 且 $\gamma(0)=p,\gamma(1)=q$.

记 $D=d(p,q)>0$. 对 $s<t$, 由恒速逼近列有 $d(\gamma_j(s),\gamma_j(t))\le L(\gamma_j)(t-s)$, 极限后为 $d(\gamma(s),\gamma(t))\le D(t-s)$. 将 $[0,1]$ 按 $0,s,t,1$ 分割, 端点三角不等式给出 $D\le Ds+D(t-s)+D(1-t)=D$. 所以每一项都必须取等, 即 $d(\gamma(s),\gamma(t))=D(t-s)$.

还需证明这个距离等距路径确实光滑, 不能直接套用仅为正则路径证明的取等结论. 固定 $a<1$, 取 $\delta>0$ 使 $\gamma([a,a+\delta])$ 留在以 $\gamma(a)$ 为中心的正常邻域. 其正常径向坐标为 $r(t)=D(t-a)$. 对 $a<s<t\le a+\delta$ 且 $t-s$ 足够小, 在 $\gamma(s)$ 的正常邻域中, 用一条光滑径向最短测地线连接 $\gamma(s),\gamma(t)$, 长度为 $D(t-s)=r(t)-r(s)$. 缩短 $t-s$ 使这段连接线留在原中心的正常极坐标环带内. 对该光滑连接线用 Gauss 引理, 长度至少为径向坐标变化; 现已取等, 因而它的角坐标不变, 两端有 $\theta(s)=\theta(t)$. 这使极限路径的角坐标在 $(a,a+\delta]$ 上局部恒定, 从而在此连通区间恒定. 故 $\gamma(t)=\Exp_{\gamma(a)}(D(t-a)e)$, 其中 $e$ 为固定单位方向.

在每个内部时刻选稍早的 $a$, 可知 $\gamma$ 局部为光滑恒速测地线; 初值唯一性使各段相容. 它的速度为 $D$, 在 $[0,1]$ 上长度为 $D$, 即全局最短. $p=q$ 时取常值路径.
\end{proof}
原课程第36--39讲给出的闭子集版本由此在二维嵌入情形补成证明. 若去掉闭性, 结论可失败: 穿孔平面中穿过孔洞的直线不能延续, 位于孔洞两侧的点间距离下确界也可能没有路径达到.
\originalsection{习题}
\begin{exercise}\label{ex:9a}
球面半径为 $R$, 点 $p,q$ 的中心夹角为 $\alpha\in[0,\pi]$. 证明 $d(p,q)=R\alpha$, 并说明最短曲线在非对跖、对跖及同点三种情况的区别.
\end{exercise}
\begin{exercise}\label{ex:9b}
在穿孔平面 $\R^2\setminus\{0\}$ 中, 证明 $p=(-1,0)$ 与 $q=(1,0)$ 的内蕴距离是2, 但不存在达到距离的路径.
\end{exercise}
\begin{exercise}\label{ex:9c}
在正常极坐标中, 证明 $A(r,\theta)=r+O(r^3)$, 并用 $K=-A_{rr}/A$ 得到 $A(r,\theta)=r-\frac16K(p)r^3+o(r^3)$. 明确 $r=0$ 的初值如何取得.
\end{exercise}

\originalchapter{平行移动与局部 Gauss--Bonnet}
\originalsection{平行移动}
把一个切向量从一点移到另一点, 不能要求它在三维空间中完全不变, 因为切平面本身改变. 内蕴的“不转动”要求其协变导数为0.
\begin{definition}
沿曲线 $\gamma$ 的切向量场 $V$ 若 $\nabla_{\dot\gamma}V=0$, 称为平行场. 给定起点向量后, 方程的解定义沿该路径的平行移动.
\end{definition}
\begin{proposition}\label{prop:parallel}
平行移动存在唯一, 保持向量内积. 在正向单位标架 $e_1,e_2$ 下, 单位向量 $V=\cos\phi\,e_1+\sin\phi\,e_2$ 平行当且仅当 $\dot\phi=-\omega(\dot\gamma)$.
\end{proposition}
\begin{proof}
坐标分量满足线性方程 $\dot v^k+\sum_{i,j}\Gamma^k_{ij}\dot x^iv^j=0$, 线性常微分方程给出存在唯一性. 度量相容性说明两个平行场的内积导数为0. 对角表达式求导, 与 $JV=-\sin\phi e_1+\cos\phi e_2$ 的分量为 $\dot\phi+\omega(\dot\gamma)$, 得最后条件.
\end{proof}
平行移动依赖路径. 对在同一坐标片中的正向闭边界 $\partial D$, 起点和终点切平面相同. 净旋转角模 $2\pi$ 称为绕该边界的 holonomy 角. 由命题\ref{prop:connection-curvature},
\begin{equation}\label{eq:holonomy}
 \Delta\phi=-\oint_{\partial D}\omega=\int_D K\dd A\pmod{2\pi}.
\end{equation}
它将曲率解释为小回路平行移动的角缺陷. 等式的实数代表依赖选取的角分支, 几何旋转只按模 $2\pi$ 定义.
\originalsection{测地曲率与 Darboux 分解}
\begin{definition}
在定向曲面上, 单位速度正则曲线的有向测地曲率为 $k_g=\langle\nabla_TT,N\times T\rangle$. 任意正则参数下,
\[
 k_g=\frac{\det(\dot\gamma,\ddot\gamma,N)}{|\dot\gamma|^3}.
\]
\end{definition}
空间加速度在正交基 $T,N\times T,N$ 下有分解
\begin{equation}\label{eq:darboux}
 \gamma''=k_g(N\times T)+k_nN,\qquad \kappa^2=k_g^2+k_n^2
\end{equation}
其中使用弧长参数, $k_n=\II(T,T)$. 第一式来自 $\gamma''\perp T$ 及前两种曲率定义, 第二式是正交分量平方和. 所以法曲率与测地曲率描述不同弯曲分量.
\begin{proposition}\label{prop:kg-angle}
若 $T=\cos\theta e_1+\sin\theta e_2$, 则 $k_g=\dd\theta/\dd s+\omega(T)$. 单位速度曲线为测地线当且仅当 $k_g=0$.
\end{proposition}
\begin{proof}
对 $T$ 使用乘积法则和 $\nabla e_1=\omega e_2,\nabla e_2=-\omega e_1$, 得 $\nabla_TT=(\theta'+\omega(T))(N\times T)$. 无切向加速度的条件就是该系数为0. 若不是恒速参数, $k_g=0$ 只说明轨迹可以弧长化为测地线, 并不消除沿切向的速度变化.
\end{proof}
\begin{example}
半径 $R$ 球面取外法向. 求北半球纬圆 $\theta=\theta_0\in(0,\pi/2)$ 沿 $\varphi$ 增加方向的测地曲率和北极帽的曲率积分.
\end{example}
\begin{solution}
记球坐标 $X=R(\sin\theta\cos\varphi,\sin\theta\sin\varphi,\cos\theta)$. 纬圆空间曲率为 $1/(R\sin\theta_0)$, 主法向指向纬圆圆心. 直接将其加速度与 $N\times T$ 内积, 得 $k_g=\cot\theta_0/R$. 弧长元为 $R\sin\theta_0\dd\varphi$, 故 $\int k_g\dd s=2\pi\cos\theta_0$. 北极帽面积为 $2\pi R^2(1-\cos\theta_0)$, 乘 $K=1/R^2$ 得 $\int K\dd A=2\pi(1-\cos\theta_0)$. 两项之和为 $2\pi$.
\end{solution}
\originalsection{带角点的圆盘公式}
本节区域为包含在一个正向坐标片内的拓扑圆盘 $D$, 边界由有限条 $C^2$ 正则弧组成, 交点处有非零夹角, 无尖点回折. 沿正向边界在第 $j$ 个角点, 从入切向到出切向的有向外转角为 $\beta_j\in(-\pi,\pi)$, 内角 $\alpha_j=\pi-\beta_j\in(0,2\pi)$. 光滑边界没有角点项.
\begin{theorem}[局部 Gauss--Bonnet]\label{thm:local-gb}
在上述条件下,
\begin{equation}\label{eq:local-gb}
 \int_DK\dd A+\int_{\partial D}k_g\dd s+\sum_j\beta_j=2\pi.
\end{equation}
\end{theorem}
\begin{proof}
坐标片上可用 Gram--Schmidt 取得全局正向单位标架. 在每条光滑边上取切向角 $\theta$, 于角点加上有向跳跃 $\beta_j$. 切向一周相对于此标架的总角变化为 $2\pi$. 为解释这点, 参数平面中边界的切向按定理\ref{thm:turning}及其角点版本总转 $2\pi$. 从坐标基换到曲面正交基的变换是 $D$ 上的正行列式矩阵场. 它在圆盘上可连续变形到常矩阵场: 正定对称部分可线性变形到恒等, 旋转部分在单连通圆盘上可取连续角函数并收缩. 此变形不使非零切向量成为零, 所以边界单位切向映射的旋转数不变. 因而相对于所选标架也总转 $2\pi$. 角点版本亦可将角点在小邻域光滑化, 外角正是被替换小弧的极限转角.

逐边积分命题\ref{prop:kg-angle}, 得
\[
 \int_{\partial D}k_g\dd s+\sum_j\beta_j
 =2\pi+\oint_{\partial D}\omega.
\]
连接形式光滑定义于整个圆盘, Green 公式与 $\dd\omega=-K\dd A$ 给出 $\oint\omega=-\int_DK\dd A$, 即所需等式.
\end{proof}
\begin{corollary}[测地三角形]\label{cor:triangle}
包含在单一坐标圆盘中的简单测地三角形, 内角为 $\alpha_1,\alpha_2,\alpha_3$, 满足
\[
 \alpha_1+\alpha_2+\alpha_3=\pi+\int_DK\dd A.
\]
\end{corollary}
\begin{proof}
三边 $k_g=0$, 三角点外角和为 $3\pi-\sum\alpha_j$, 代入局部公式. 本结论对跨多个坐标片的三角形会由整体带边界版本推广.
\end{proof}
\begin{center}
\begin{tikzpicture}[scale=1.05,>=Latex]
\draw[thick,structurecolor] (0,0) to[bend right=12] (3.6,0) to[bend right=12] (1.7,2.5) to[bend right=12] (0,0);
\node[below left] at (0,0) {$A$};\node[below right] at (3.6,0) {$B$};\node[above] at (1.7,2.5) {$C$};
\node at (.55,.45) {$\alpha_1$};\node at (3.0,.45) {$\alpha_2$};\node at (1.7,1.95) {$\alpha_3$};
\draw[->] (1.2,-.28)--(2.2,-.28);
\node at (1.75,.95) {$\displaystyle\int_D K\dd A$};
\end{tikzpicture}
\end{center}
图为曲面圆盘上的三角形示意, 不是在参数平面中测量三角形欧氏角. 曲面内角应由第一基本形式计算.
\originalsection{习题}
\begin{exercise}\label{ex:10a}
在定向曲面上反转法向 $N$, 但保持同一参数方向的曲线, 说明 $k_g,K,k_n$ 如何变化. 再说明 Gauss--Bonnet 中为何必须同时调整正向边界方向.
\end{exercise}
\begin{exercise}\label{ex:10b}
半径 $R$ 的球面上, 三条互相垂直的过心平面围成一个球面八分之一三角形. 求三个内角、面积、曲率积分, 验证三角形公式.
\end{exercise}
\begin{exercise}\label{ex:10c}
若单位向量沿球面北极帽边界平行移动一周, 北极帽的极角半径为 $\theta_0\in(0,\pi)$, 求净旋转角模 $2\pi$. 当 $\theta_0=\pi/2$ 时为何结果不与球面非零曲率矛盾?
\end{exercise}

\originalchapter{整体 Gauss--Bonnet 与 Gauss 映射}
\originalsection{有限三角剖分与曲率总量}
本章只引用拓扑中的两个工具: 紧致光滑曲面区域有有限光滑三角剖分, 可细分到每个三角形包含在一个坐标片内; 对任一此类剖分, $\chi=V-E+F$ 为 Euler 示性数, 不依赖剖分. 边界是有限条光滑闭曲线或分段光滑曲线, 三角剖分与其角点相容. 这些是拓扑工具, 不以曲率公式证明其存在与不变性.
\begin{theorem}[Gauss--Bonnet]\label{thm:global-gb}
设 $M$ 为紧致定向光滑曲面区域, 边界分段 $C^2$ 正则, 角点内角属于 $(0,2\pi)$, 且边界无自交. 边界按区域在左定向, 角点外角为 $\beta_j$. 则
\begin{equation}\label{eq:global-gb}
 \int_MK\dd A+\int_{\partial M}k_g\dd s+\sum_j\beta_j=2\pi\chi(M).
\end{equation}
无边界时即 $\int_MK\dd A=2\pi\chi(M)$.
\end{theorem}
\begin{proof}
取上述有限细三角剖分, 对每个三角形应用定理\ref{thm:local-gb}. 内部公共边由两个三角形沿相反方向走过, 测地曲率变号, 因而边积分相消; 边界边留下原边界积分. 记内部顶点数为 $V_i$, 边界顶点数为 $V_b$, 三角形数为 $F$. 每个三角形有3个角, 各外角为 $\pi$ 减该三角形内角, 所以把各局部公式相加得
\[
 \int_MK\dd A+\int_{\partial M}k_g\dd s
 =\sum_{\text{所有三角形角}}\alpha-\pi F.
\]
内部顶点周围各内角和为 $2\pi$. 光滑边界顶点周围为 $\pi$, 原有角点处则为其区域内角 $\alpha_j$. 在两侧加上原边界外角 $\pi-\alpha_j$, 得右侧总计 $2\pi V_i+\pi V_b-\pi F$.

记内部边数 $E_i$, 边界边数 $E_b$. 每条闭边界的边数等于顶点数, 故 $E_b=V_b$, 且三角形边计数给出 $3F=2E_i+E_b$. 因而
\[
 2\pi\chi=2\pi(V_i+V_b-E_i-E_b+F)
 =2\pi V_i+\pi V_b-\pi F,
\]
正是前面所得右侧. 无边界时取 $V_b=E_b=0$ 即得相应版本.
\end{proof}
\begin{example}
计算球面、环面以及有 $b$ 条边界的亏格 $g$ 连通定向紧致曲面区域的曲率约束.
\end{example}
\begin{solution}
球面 $\chi=2$, 所以总曲率为 $4\pi$; 半径改变使 $K$ 乘 $R^{-2}$、面积乘 $R^2$, 积分不变. 环面 $\chi=0$, 总曲率为0, 与第6章直接积分一致.

拓扑给出一般 $\chi=2-2g-b$. 若边界光滑且全为测地线, $k_g=0$, 则 $\int_MK\dd A=2\pi(2-2g-b)$. 例如带三个测地边界的亏格0区域, 即裤形区域, 总曲率为 $-2\pi$. 这说明在此条件下不可能处处 $K\ge0$.
\end{solution}
\begin{corollary}
紧致连通无边界定向曲面若 $K>0$, 则其亏格为0. 若 $K\le0$ 且不恒为0, 则亏格至少为2. 若亏格为1且 $K\ge0$ 或 $K\le0$, 则 $K\equiv0$.
\end{corollary}
\begin{proof}
用 $\int K=4\pi(1-g)$ 及连续函数积分的符号. 最后一项总积分为0, 不变号连续函数若不恒零会在某个开集上严格同号, 积分不为0, 矛盾.
\end{proof}
这里的亏格分类来自拓扑, 而曲率负责限制分类中哪些情况能实现. 对嵌入 $\R^3$ 的紧致曲面, 下节还会排除处处平直的可能.
\originalsection{Gauss 映射与椭圆点}
\begin{definition}
定向曲面的 Gauss 映射为 $N:M\to S^2$, 送每点到其单位法向. 目标球面按外法向定向.
\end{definition}
\begin{proposition}\label{prop:gauss-jacobian}
Gauss 映射的有向面积 Jacobian 为 $K$, 即 $N^*\dd A_{S^2}=K\dd A_M$. 绝对面积缩放为 $|K|$.
\end{proposition}
\begin{proof}
$p$ 与 $N(p)$ 的切平面均为 $N(p)^\perp$, 而正向切基以与 $N(p)$ 的叉积相容为条件, 两边可用同一正向正交基. $\dd N=-S$ 在二维上行列式为 $\det(-S)=\det S=K$. 面积换元的绝对值则为 $|K|$.
\end{proof}
这个等式解释了 $K$ 的“法向球面积密度”含义, 但在全局非单射时要计覆盖重数. 总有向面积有正负抵消, 不能直接等同于 Gauss 像点集的普通面积.
\begin{proposition}\label{prop:elliptic-point}
非空紧致无边界嵌入曲面至少有一点 $K>0$. 取原点在曲面之外并令 $|p|=R$ 在曲面达到最大, 在此处适当选外径向法向, 两主曲率都不大于 $-1/R$.
\end{proposition}
\begin{proof}
在最大点 $p$, 函数 $|X|^2$ 一阶导数为0, 所以 $p$ 垂直切平面, 可取 $N(p)=p/R$. 对任意单位切向 $V$, 取曲面曲线 $\gamma$ 经过 $p$ 且速度为 $V$. 二阶导数非正, 得
\[
 0\ge\tfrac12(|\gamma|^2)''=1+\langle p,\gamma''\rangle
 =1+R\II(V,V).
\]
因此所有单位方向的法曲率不大于 $-1/R$, 特别两个主曲率都严格为负, $K\ge1/R^2>0$. 反转法向仍保持 $K>0$.
\end{proof}
\begin{corollary}
任何紧致无边界嵌入环面都有 $K>0$ 的点和 $K<0$ 的点.
\end{corollary}
\begin{proof}
正值点由上命题得到. 总曲率为0, 连续性使正曲率在一个开集上贡献正积分, 故别处必须有负值.
\end{proof}
\originalsection{正曲率与整体凸性}
这里使用覆盖理论中的事实: 紧致连通空间到连通流形的局部微分同胚是有限覆盖; 连通覆盖的目标若为单连通球面, 则覆盖为一层. 在本册二维情形可直接说明第一点: 原像紧致而离散, 所以有限; 由逆函数定理在各原像点取互不相交小邻域, 再借紧致性排除其外的新原像, 得均匀覆盖邻域. 第二点是明确引用的拓扑结论.
\begin{theorem}[二维 Hadamard 凸性]\label{thm:hadamard}
紧致连通无边界嵌入曲面若 $K$ 处处非零, 则 $K>0$, Gauss 映射为微分同胚, 且每个切平面都是支撑平面. 它围成一个严格凸体.
\end{theorem}
\begin{proof}
由连续性和连通性, $K$ 不变号; 命题\ref{prop:elliptic-point}强迫其处处为正. 由于 $K>0$, 每点两种法向中恰有一种令形算子负定; 这个选择在局部光滑且在重叠片一致, 给出全局光滑法向. 取此法向, $\dd N=-S$ 处处可逆. Gauss 映射像是开集, 又因紧致为闭集, 所以覆盖整个 $S^2$. 由上述覆盖与单连通性, 它是一层覆盖, 即微分同胚.

主曲率乘积处处为正, 两者同号且不会越过0. 适当反转整体法向使最大距离点的形算子负定, 连通性使它处处负定. 对任意方向 $a\in S^2$, 高度函数 $h_a(p)=\langle p,a\rangle$ 的临界点恰为 $N(p)=a$ 或 $N(p)=-a$, 各有唯一一点. 在 $N=a$ 点, 高度 Hessian 为 $\II$, 负定, 是局部最大; 在 $N=-a$ 点它为 $-\II$, 正定, 是局部最小. 紧致性保证全局最大最小存在, 所以这两个点也是唯一全局最大和最小. 因而所有切平面都是支撑平面, 各方向最大支撑点唯一.

令 $C$ 为曲面的凸包. 它紧致: 三维空间中的凸组合总能缩到至多四个点, 因为超过四个点时仿射线性相关, 可沿相关系数调整权重直到至少一个权重为0; 重复即可. 因而凸包是紧集 $M^4$ 与紧致权重单纯形之连续像. 曲面不共面, 所以凸包有内部. 每个 $\partial C$ 点都有支撑平面: 取外部点趋近该边界点, 外部点到紧凸集的最近点给出支撑法向, 单位法向取收敛子列得到极限支撑面. 最近点性质的证明是对任意 $y\in C$ 比较到线段 $x+t(y-x)$ 的距离平方在 $t=0$ 的右导数, 得 $\langle a,y-x\rangle\le0$.

最大支撑值由曲面取得. 若凸组合落在对应支撑面上, 其所有正权重分量都必须达到该高度最大值, 而最大点唯一, 所以支撑面与 $C$ 的交仅有该点. 每个 $\partial C$ 点因此属于曲面. 反之每个曲面点都有支撑平面, 属于 $\partial C$, 故 $M=\partial C$. 其边界不含线段: 若含线段, 线段内点的支撑平面将同时包含两端点, 与单点支撑面矛盾. 所以 $C$ 是严格凸体.
\end{proof}
该定理关于紧致嵌入曲面. 不能把它用于开曲面或任意非单射浸入, 也不能由局部 $K>0$ 推断整曲面凸.
\originalsection{习题}
\begin{exercise}\label{ex:11a}
设紧致连通无边界定向曲面处处 $K=-1$, 亏格为 $g$. 求面积, 判断哪些 $g$ 从 Gauss--Bonnet 看可能. 说明该结论是否自动给出三维欧氏空间中的完整嵌入.
\end{exercise}
\begin{exercise}\label{ex:11b}
半径 $R$ 的球面上去掉两个互不相交极帽, 剩余纬带极角为 $\theta\in[\theta_1,\theta_2]$, $0<\theta_1<\theta_2<\pi$. 计算曲率积分和两条正向边界的测地曲率积分, 核对 Euler 示性数.
\end{exercise}
\begin{exercise}\label{ex:11c}
证明平面中光滑环形区域的两条正向边界测地曲率积分之和为0, 并解释为什么分别为 $2\pi$ 和 $-2\pi$.
\end{exercise}

\originalchapter{双曲平面与常曲率模型}
\originalsection{Poincar\'e 圆盘度量}
内蕴几何允许先给正定度量, 再谈长度、曲率与测地线, 而不要求已经找到三维欧氏曲面. 这是原课程 Minkowski 双曲模型所揭示的观点. 本章采用等价的二维圆盘度量, 不把它误认成欧氏圆盘的通常长度. 对任意二维光滑正定度量, 用第7章联络定义 $K=\langle R(e_1,e_2)e_2,e_1\rangle$, 此值与单位正交标架无关: 二维曲率算子的反对称性与度量相容性使它由这一标量唯一决定, 正交换基的两个面积因子相乘为1. 正向标架仍满足 $\dd\omega=-K\dd A$. 这一定义在欧氏曲面上由 Gauss 方程与 $\det S$ 一致, 在未嵌入的圆盘上则无须形算子.
\begin{definition}
单位圆盘 $\mathbb D=\{z=u+iv\in\C:|z|<1\}$ 上的双曲度量为
\[
 \I=\frac{4(\dd u^2+\dd v^2)}{(1-u^2-v^2)^2}.
\]
其面积元为 $\dd A=4\dd u\dd v/(1-|z|^2)^2$.
\end{definition}
\begin{proposition}\label{prop:hyperbolic-K}
圆盘度量的 Gauss 曲率恒为 $-1$. 欧氏边界圆不是度量空间中的点集.
\end{proposition}
\begin{proof}
写 $f=\log2-\log(1-u^2-v^2)$, 则度量为 $e^{2f}(\dd u^2+\dd v^2)$. 有 $\Delta f=4/(1-u^2-v^2)^2$, 由命题\ref{prop:metric-K}, $K=-e^{-2f}\Delta f=-1$. 当 $|z|\to1$, 系数发散, 只在开圆盘中定义光滑正定度量; 边界不能当作普通曲面边界代入距离公式.
\end{proof}
\begin{proposition}\label{prop:disk-isometry}
对 $a\in\mathbb D$, 映射 $\phi_a(z)=(z-a)/(1-\overline a z)$ 为圆盘的等距微分同胚, 正向旋转亦为等距.
\end{proposition}
\begin{proof}
直接展开得
\[
 1-|\phi_a(z)|^2=\frac{(1-|a|^2)(1-|z|^2)}{|1-\overline a z|^2},\qquad
 |\phi_a'(z)|=\frac{1-|a|^2}{|1-\overline a z|^2}.
\]
故 $2|\phi_a'(z)||\dd z|/(1-|\phi_a(z)|^2)=2|\dd z|/(1-|z|^2)$, 长度元不变. 分母在圆盘非零, 逆映射为 $w\mapsto(w+a)/(1+\overline a w)$, 仍送圆盘到圆盘, 因而为微分同胚. 旋转对模长与欧氏长度不变, 同样成立.
\end{proof}
\originalsection{距离与测地线}
\begin{theorem}\label{thm:hyperbolic-distance}
双曲距离为
\begin{equation}\label{eq:disk-distance}
 d(z,w)=2\arctanh\left|\frac{w-z}{1-\overline z w}\right|
 =\log\frac{1+q}{1-q},\qquad q=\left|\frac{w-z}{1-\overline z w}\right|.
\end{equation}
其中 $0\le q<1$. 连接两点的最短路径轨迹唯一.
\end{theorem}
\begin{proof}
先求从0到半径 $r$ 的距离. 任意路径以极坐标表示时, 长度元满足
\[
 \frac{2\sqrt{\dot r^2+r^2\dot\theta^2}}{1-r^2}\ge\frac{2|\dot r|}{1-r^2}.
\]
积分至少为 $\int_0^r2\dd t/(1-t^2)=2\arctanh r$, 径向线达到. 对可穿过0的路径可分段或在小圆外取极坐标再令半径趋于0. 等长要求角方向不变且半径不回走, 故轨迹唯一. 用等距 $\phi_z$ 把起点送到0, 终点模长为 $q$, 得公式及唯一性. $q<1$ 来自上一命题的圆盘映射恒等式.
\end{proof}
径向路径采用双曲弧长 $s$ 时为 $z(s)=\tanh(s/2)e^{i\theta_0}$. 最短路径的一阶能量变分为0, 故它是测地线; 等距保持能量变分与测地线. 初值唯一性因此说明所有非恒定测地线轨迹都是径向直径的圆盘等距像.
\begin{corollary}\label{cor:disk-geodesics}
非恒定双曲测地线的最大延拓轨迹在欧氏图中为直径, 或与单位圆正交的圆盘内圆弧. 任意两个不同点位于唯一这样的轨迹上. 限制参数区间只截取该轨迹的一段.
\end{corollary}
\begin{proof}
分式线性变换把广义圆, 即圆或直线, 送到广义圆. 这一事实可由将 $\phi_a^{-1}$ 代入广义圆的方程 $A|z|^2+Bz+\overline B\overline z+C=0$ 后清除分母验证. 它还保角, 因复导数非零在每点是伸缩与旋转. 所以直径与边界的直角经映射后仍为直角. 上一定理的径向唯一性给出两点间唯一弧, 初值存在唯一性排除其他轨迹.
\end{proof}
\begin{center}
\begin{tikzpicture}[scale=2]
\draw[thin] (0,0) circle (1);
\draw[thick,structurecolor] (-1,0)--(1,0);
\draw[thick,structurecolor,domain=150:210,samples=60] plot ({2+sqrt(3)*cos(\x)},{sqrt(3)*sin(\x)});
\fill (.5,.866025) circle (.7pt);\fill (.5,-.866025) circle (.7pt);
\node[anchor=west] at (1.12,.55) {边界不属于 $\mathbb D$};
\node[anchor=east] at (-1.1,-.5) {直径};
\node[anchor=west] at (.7,-1.1) {正交圆弧};
\end{tikzpicture}
\end{center}
\begin{proposition}
双曲圆盘在此距离下完备. 半径 $\rho$ 的双曲圆周长为 $2\pi\sinh\rho$, 圆盘面积为 $2\pi(\cosh\rho-1)$.
\end{proposition}
\begin{proof}
若点列为双曲 Cauchy, 到0的距离有界, 距离公式使欧氏模长统一小于某个 $r_0<1$. 又度量系数至少为4, 故它是欧氏 Cauchy, 收敛到圆盘内部点; 在紧圆盘内度量与欧氏度量一致有界比较, 因而双曲收敛. 这证明完备, 并解释边界为什么在无穷远.

令 $r=\tanh(\rho/2)$, 则极坐标度量化为 $\dd\rho^2+\sinh^2\rho\dd\theta^2$: 径向项来自 $\dd\rho=2\dd r/(1-r^2)$, 环向系数来自 $2r/(1-r^2)=\sinh\rho$. 圆周积分给出 $2\pi\sinh\rho$, 面积积分 $\int_0^{2\pi}\int_0^\rho\sinh t\dd t\dd\theta$ 给出所示面积.
\end{proof}
\originalsection{角亏与局部欧氏实现}
\begin{proposition}
非退化有限双曲测地三角形的面积为 $\pi-(\alpha+\beta+\gamma)$, 因而内角和小于 $\pi$. 有 $n$ 条测地边的简单凸多边形面积为 $(n-2)\pi-\sum_j\alpha_j$.
\end{proposition}
\begin{proof}
此类区域为紧致圆盘, 边界测地曲率为0. 带角点 Gauss--Bonnet 的推导只依赖度量联络和 $\dd\omega=-K\dd A$, 因而对圆盘内蕴度量同样适用. 代入 $K=-1$ 与外角 $\pi-\alpha_j$ 得公式. 非退化三角形有正面积, 所以角和严格小于 $\pi$.
\end{proof}
\begin{example}
在三维欧氏空间局部构造 $K=-1$ 的旋转曲面, 并判断它是否给出了完整双曲圆盘的全局实现.
\end{example}
\begin{solution}
取母线弧长 $s<0$, $r(s)=e^s$, $z'(s)=\sqrt{1-e^{2s}}$, 可取任一积分常数. 由于 $r'^2+z'^2=1$, 第6章公式给出 $K=-r''/r=-1$. 这是伪球面的一个正则部分, 经向长度到 $s=0$ 有限. 在该边界 $z'$ 趋于0, 母线二阶导数出现发散, 不能把当前参数域平滑补成完整双曲平面. 有限端的缺失已足以说明此片不完备, 因而不是完整圆盘的等距全局模型. 本书不以此例证明一般负曲率完整曲面的全局不可嵌入定理.
\end{solution}
常曲率模型展示了同样的几何关系: 球面小圆周长小于欧氏 $2\pi\rho$, 双曲平面大于它; 测地三角形角和相应大于或小于 $\pi$. 这些结论在各自公式中有明确适用半径与区域, 不将示意图代替证明.
\originalsection{习题}
\begin{exercise}\label{ex:12a}
在双曲圆盘上求 $d(0,1/2)$、$d(-1/2,1/2)$, 并给出沿实轴的单位速度参数.
\end{exercise}
\begin{exercise}\label{ex:12b}
双曲平面一个测地三角形内角为 $\pi/3,\pi/4,\pi/6$. 求面积. 与单位球面上具有同样三个内角的测地圆盘三角形比较, 判断后者是否可能.
\end{exercise}
\begin{exercise}\label{ex:12c}
求度量 $\I=(\dd x^2+\dd y^2)/y^2$, $y>0$, 的曲率, 并证明竖直直线为测地线轨迹, 求其恒速参数及两点 $(x_0,y_1),(x_0,y_2)$ 的距离.
\end{exercise}

\par\Needspace{13\baselineskip}
\originalchapter{习题解答}
本章与前面各章的习题编号用链接相互对应. 各题为本书根据所讲概念另设或整理的标准训练题, 不宣称来自未核实的 MIT 作业题号.
\par\Needspace{9\baselineskip}
\originalsection{弧长与平面曲率}
\par\Needspace{5\baselineskip}
\noindent\textbf{习题\ref{ex:1a}.}
\begin{solution}
微分得 $\gamma'=e^t(\cos t-\sin t,\sin t+\cos t)$, $\gamma''=2e^t(-\sin t,\cos t)$. 所以 $v=\sqrt2e^t$, 行列式为 $2e^{2t}$, 代入曲率公式得 $k=e^{-t}/\sqrt2$. 在 $t=0$ 有 $\gamma=(1,0)$, $T=(1,1)/\sqrt2$, $n=(-1,1)/\sqrt2$, $k=1/\sqrt2$. 圆心为 $\gamma+n/k=(0,1)$, 半径为 $1/k=\sqrt2$.
\end{solution}
\par\Needspace{5\baselineskip}
\noindent\textbf{习题\ref{ex:1b}.}
\begin{solution}
利用 $n'=-kT$, 得
\[
 C'=T+\frac{n'}k-\frac{k'}{k^2}n=-\frac{k'}{k^2}n.
\]
因 $k\ne0$, 圆心轨迹正则当且仅当 $k'\ne0$. 在曲率临界点它的速度为0, 所给参数下失去正则性; 这不自动保证该点是尖点, 还需更高阶导数与非退化条件. 圆的 $k$ 恒定, 圆心轨迹是一点, 是最直接的例外.
\end{solution}
\par\Needspace{5\baselineskip}
\noindent\textbf{习题\ref{ex:1c}.}
\begin{solution}
取 $F=x^2+2y^2$, 梯度 $(2x,4y)$, Hessian 为 $\operatorname{diag}(2,4)$. 在 $(1,0)$, $J\nabla F=(0,2)$, Hessian 二次型为16, 梯度长度三次方为8, 得 $k=2$. 在 $(0,1/\sqrt2)$, $J\nabla F=(-2\sqrt2,0)$, 二次型为16, 分母为 $16\sqrt2$, 得 $k=1/\sqrt2$. 椭圆半轴 $a=1,b=1/\sqrt2$, 前章公式给出 $a/b^2=2$ 与 $b/a^2=1/\sqrt2$, 一致.
\end{solution}
\par\Needspace{9\baselineskip}
\originalsection{闭曲线}
\par\Needspace{5\baselineskip}
\noindent\textbf{习题\ref{ex:2a}.}
\begin{solution}
相对于指定周期, 它不能简单, 因简单正向曲线的总有向曲率只能为 $2\pi$, 反向为 $-2\pi$. 例子取 $\alpha(s)=R(\cos(s/R),\sin(s/R))$, 指定周期 $L=4\pi R$. 它为单位速度, $k=1/R>0$, 此周期走两圈, 总曲率为 $4\pi$. 该映射也有较小周期 $2\pi R$, 若用较小周期定义闭曲线, 则它是简单圆. 这说明必须固定所研究的一周.
\end{solution}
\par\Needspace{5\baselineskip}
\noindent\textbf{习题\ref{ex:2b}.}
\begin{solution}
令 $n_o=(\cos\varphi,\sin\varphi)$, $t=Jn_o$, 用 $\alpha=hn_o+h't$ 定义. 有 $h'=-2b\sin2\varphi$, $h''=-4b\cos2\varphi$, 所以
\[
 \rho=h+h''=a-3b\cos2\varphi>0,\qquad
 k=\frac1{a-3b\cos2\varphi}.
\]
曲线闭合且 $\alpha'=\rho t\ne0$. 为不循环地证明简单, 固定 $\varphi_0$ 并考察高度 $f(\varphi)=\langle\alpha(\varphi),n_o(\varphi_0)\rangle$. 其导数为 $\rho\sin(\varphi_0-\varphi)$, 在 $(\varphi_0,\varphi_0+\pi)$ 为负, 在下一半周为正. 因而一周内唯一最大点是 $\varphi_0$. 若另一个参数对应同一点, 它也达到该高度最大值, 矛盾. 于是曲线简单.

由 $\dd s=\rho\dd\varphi$, 有 $\dd k/\dd s=-6b\sin2\varphi/\rho^3$. $b\ne0$, 临界参数恰为 $0,\pi/2,\pi,3\pi/2$, 正好四个.
\end{solution}
\par\Needspace{5\baselineskip}
\noindent\textbf{习题\ref{ex:2c}.}
\begin{solution}
周长为 $\int\rho\dd\varphi=\int(h+h'')\dd\varphi=\int h\dd\varphi$, 最后一项由周期性 $\int h''=0$. 面积由 Green 公式为 $\frac12\int\det(\alpha,\alpha')\dd\varphi$. 因 $\alpha=hn_o+h't$, $\alpha'=\rho t$, $\det(n_o,t)=1$, 得行列式 $h\rho$. 因而
\[
 A=\tfrac12\int h(h+h'')\dd\varphi
 =\tfrac12\int(h^2-h'^2)\dd\varphi,
\]
分部积分的边界项因周期而为0. 曲线正向简单使此有向面积等于区域通常面积.
\end{solution}
\par\Needspace{9\baselineskip}
\originalsection{空间曲线}
\par\Needspace{5\baselineskip}
\noindent\textbf{习题\ref{ex:3a}.}
\begin{solution}
$\gamma'=(1,2t,3t^2)$, $\gamma''=(0,2,6t)$, $\gamma'''=(0,0,6)$. 叉积为 $(6t^2,-6t,2)$, 平方长度为 $4(1+9t^2+9t^4)$, 三重积为12. 因而
\[
 \kappa=\frac{2\sqrt{1+9t^2+9t^4}}{(1+4t^2+9t^4)^{3/2}},\qquad
 \tau=\frac3{1+9t^2+9t^4}.
\]
曲率处处正, 挠率处处正, 由平面性判据, 任意非空开区间都不可能落在一个固定平面上.
\end{solution}
\par\Needspace{5\baselineskip}
\noindent\textbf{习题\ref{ex:3b}.}
\begin{solution}
记 $r=\alpha-p$. $|r|^2=R^2$ 一次微分给 $\langle r,T\rangle=0$, 二次微分给 $1+\kappa\langle r,n\rangle=0$, 即 $\langle r,n\rangle=-1/\kappa$. 因而 $r=-n/\kappa+cb$. 再微分 $\langle r,n\rangle=-1/\kappa$, 得 $\langle r,n'\rangle=\kappa'/\kappa^2$. Frenet 方程使左侧为 $\tau c$, 所以在 $\tau\ne0$ 区间
\[
 r=-\frac1\kappa n+\frac{\kappa'}{\kappa^2\tau}b,\qquad
 R^2=\frac1{\kappa^2}+\left(\frac{\kappa'}{\kappa^2\tau}\right)^2.
\]
这是位于给定球面时的必要关系. 单独只知右侧恒定, 不能在此解答中未经进一步论证宣称已构造固定球心; 原题只要求从已在球面上推出该式.
\end{solution}
\par\Needspace{5\baselineskip}
\noindent\textbf{习题\ref{ex:3c}.}
\begin{solution}
设投影 $\beta=\alpha-\langle\alpha,U\rangle U$. 有 $\langle T,U\rangle=d/\sqrt{1+d^2}$, $U$ 恒定, 因而 $\beta'=T-dU/\sqrt{1+d^2}$, 速度 $v_\beta=1/\sqrt{1+d^2}>0$. 又 $\beta''=\kappa n$, 且 $n\perp U$, $n\perp\beta'$. 因速度为常数, 投影曲率为 $|\beta''|/v_\beta^2=(1+d^2)\kappa$. 它虽在平面上, 仍不必是圆, 因为 $\kappa$ 可随弧长变化.
\end{solution}
\par\Needspace{9\baselineskip}
\originalsection{曲面度量}
\par\Needspace{5\baselineskip}
\noindent\textbf{习题\ref{ex:4a}.}
\begin{solution}
$X_u=(1,0,f_u)$, $X_v=(0,1,f_v)$, 所以
\[
 \I=(1+f_u^2)\dd u^2+2f_uf_v\dd u\dd v+(1+f_v^2)\dd v^2,\qquad
 \dd A=\sqrt{1+f_u^2+f_v^2}\dd u\dd v.
\]
对 $f=(u^2+v^2)/2$, 极坐标下 $\dd A=\sqrt{1+r^2}r\dd r\dd\theta$. 面积为
\[
 2\pi\int_0^a r\sqrt{1+r^2}\dd r
 =\frac{2\pi}3\bigl((1+a^2)^{3/2}-1\bigr).
\]
积分区域在参数平面, 被积面积因子反映倾斜程度.
\end{solution}
\par\Needspace{5\baselineskip}
\noindent\textbf{习题\ref{ex:4b}.}
\begin{solution}
记 $q=u^2+v^2$, 分子范数平方为 $4q+(q-1)^2=(q+1)^2$, 故 $|X|=1$. 它不会到北极 $(0,0,1)$, 反之球面非北极点的逆为 $(u,v)=(x/(1-z),y/(1-z))$, 所以为单射光滑参数.

用商法则微分后可求得 $\langle X_u,X_u\rangle=\langle X_v,X_v\rangle=4/(1+q)^2$, $\langle X_u,X_v\rangle=0$. 于是 $\I=4(\dd u^2+\dd v^2)/(1+q)^2$, 正定且与欧氏度量成同一正因子, 所以保角. 它不保长度, 因伸缩因子随点变化.
\end{solution}
\par\Needspace{5\baselineskip}
\noindent\textbf{习题\ref{ex:4c}.}
\begin{solution}
有 $\I=(1+a^2)\dd r^2+r^2\dd\theta^2$. 令 $\rho=\sqrt{1+a^2}\,r$, $\psi=\theta/\sqrt{1+a^2}$, 则 $\I=\dd\rho^2+\rho^2\dd\psi^2$, 是平面极坐标度量. 在不绕过切缝的小邻域, $(\rho,\psi)\mapsto(\rho\cos\psi,\rho\sin\psi)$ 给出局部等距. 锥上一周 $\theta$ 增加 $2\pi$, 展开扇形角为 $2\pi/\sqrt{1+a^2}$. 顶点不在定义域, 局部可展并不能恢复顶点正则性.
\end{solution}
\par\Needspace{9\baselineskip}
\originalsection{曲面曲率}
\par\Needspace{5\baselineskip}
\noindent\textbf{习题\ref{ex:5a}.}
\begin{solution}
令 $r^2=x^2+y^2$, $D=1+r^2/a^2$. 图像的 Hessian 为 $a^{-1}I_2$, 所以 $K=1/(a^2D^2)$. 用平均曲率公式, 得
\[
 H=\frac{2+r^2/a^2}{2aD^{3/2}}.
\]
在 $r>0$ 点, 径向和环向为主方向, 相应主曲率为 $1/(aD^{3/2})$ 和 $1/(a\sqrt D)$. 这可用径向单位参数方向的第一基本形式值 $D$ 与第二基本形式值 $1/(a\sqrt D)$ 相除, 环向则第一基本形式值1而第二值相同. 两者处处正, 所有点为椭圆点. 只有 $D=1$, 即原点, 两主曲率相等, 是唯一脐点; 原点值都为 $1/a$.
\end{solution}
\par\Needspace{5\baselineskip}
\noindent\textbf{习题\ref{ex:5b}.}
\begin{solution}
Euler 公式给 $k_n=2\cos^2\theta-\sin^2\theta=3\cos^2\theta-1$. 零曲率方向满足 $\cos^2\theta=1/3$, 即 $\tan^2\theta=2$, 是两条无向直线. 最大值为2, 在第一主方向; 最小值为 $-1$, 在第二主方向. 对所有角作平均, $\cos^2$ 与 $\sin^2$ 的均值均为 $1/2$, 所以平均法曲率为 $(2-1)/2=1/2=H$.
\end{solution}
\par\Needspace{5\baselineskip}
\noindent\textbf{习题\ref{ex:5c}.}
\begin{solution}
对 $N=\nabla F/|\nabla F|$ 求导, 与切向 $W$ 取内积时归一化因子的导数项沿 $\nabla F$, 被切向正交消去. 所以 $\langle\dd N(V),W\rangle=\operatorname{Hess}F(V,W)/|\nabla F|$, 取负得到 $\II$.

在椭球点 $(a,0,0)$, 外法向为 $(1,0,0)$, 切平面由 $e_y,e_z$ 正交张成. $|\nabla F|=2/a$, Hessian 为 $\operatorname{diag}(2/a^2,2/b^2,2/c^2)$. 限制切平面得主曲率 $-a/b^2$ 与 $-a/c^2$, 分别沿 $y,z$ 方向; $K=a^2/(b^2c^2)>0$. 若改内法向, 两者同变号.
\end{solution}
\par\Needspace{9\baselineskip}
\originalsection{旋转与极小曲面}
\par\Needspace{5\baselineskip}
\noindent\textbf{习题\ref{ex:6a}.}
\begin{solution}
管外法向下 $k_\varphi=-1/a$, $k_\theta=-\cos\varphi/(R+a\cos\varphi)$, 故
\[
 H=-\frac{R+2a\cos\varphi}{2a(R+a\cos\varphi)}.
\]
若 $R=2a$, 令 $H=0$ 得 $\cos\varphi=-1$, 即内赤道 $\varphi=\pi$. 其余点 $H\ne0$, 所以环面不为极小曲面. 某一条曲线上平均曲率为0只是点集上的现象.
\end{solution}
\par\Needspace{5\baselineskip}
\noindent\textbf{习题\ref{ex:6b}.}
\begin{solution}
$X_r\times X_\theta=r(-a\cos\theta,-a\sin\theta,1)$, 法向为此向量除以 $r\sqrt{1+a^2}$. 第一形式系数 $E=1+a^2,F=0,G=r^2$; 第二形式 $e=f=0,g_2=ar/\sqrt{1+a^2}$. 故径向主曲率为0, 环向为 $a/(r\sqrt{1+a^2})$, 平均曲率为 $a/(2r\sqrt{1+a^2})$, 而 $K=0$. 这说明零 Gauss 曲率不等于极小.
\end{solution}
\par\Needspace{5\baselineskip}
\noindent\textbf{习题\ref{ex:6c}.}
\begin{solution}
取原点在曲面之外, 由紧致性取 $|p|=R>0$ 最大点. 在该点所有切向均垂直 $p$, 选法向 $N=p/R$. 对单位切向 $V$ 取通过点的曲面曲线, 距离平方二阶导数非正给出 $1+R\II(V,V)\le0$. 两个主方向上因此均有 $k_i\le-1/R$, 故 $H\le-1/R<0$. 若曲面极小则 $H=0$, 矛盾. 此论证不需要先选全局法向, 因 $H=0$ 对法向反转不变.
\end{solution}
\par\Needspace{9\baselineskip}
\originalsection{内蕴联络}
\par\Needspace{5\baselineskip}
\noindent\textbf{习题\ref{ex:7a}.}
\begin{solution}
度量矩阵为 $\operatorname{diag}(1,r^2)$, 只对 $u$ 变化. 代入 Christoffel 公式得
\[
 \Gamma^u_{vv}=-rr',\qquad \Gamma^v_{uv}=\Gamma^v_{vu}=r'/r,
\]
其余为0. 正交标架 $e_1=\partial_u,e_2=r^{-1}\partial_v$ 的连接形式为 $\omega=r'\dd v$, 外微分为 $r''\dd u\wedge\dd v$. 由 $\dd\omega=-K r\dd u\wedge\dd v$ 得 $K=-r''/r$. 同一个公式适用于不先给旋转嵌入的抽象度量.
\end{solution}
\par\Needspace{5\baselineskip}
\noindent\textbf{习题\ref{ex:7b}.}
\begin{solution}
$f_u=2u,f_v=2v$, 故 $\omega=-2v\dd u+2u\dd v$. 其外微分为 $4\dd u\wedge\dd v$, 面积元为 $e^{2(u^2+v^2)}\dd u\dd v$, 所以 $K=-4e^{-2(u^2+v^2)}<0$. 平面 Gauss 曲率为0, 局部等距必须保持曲率, 因而这个度量任何点附近都不能等距于欧氏平面.
\end{solution}
\par\Needspace{5\baselineskip}
\noindent\textbf{习题\ref{ex:7c}.}
\begin{solution}
取原正向单位标架及连接形式 $\omega$. 因 $K=0$, $\dd\omega=0$. 在足够小的星形坐标邻域, 闭一形式可积分为一个函数: 若 $\omega=A\dd u+B\dd v$, 从固定底边先横后竖积分, 用 $B_u=A_v$ 核查两个偏导即得到势函数. 取 $\psi$ 满足 $\dd\psi=-\omega$, 将标架旋转角 $\psi$, 新连接形式为0, 所以两标架向量 $\widetilde e_i$ 对所有方向都平行.

定义其对偶一形式 $\eta_i(V)=\langle\widetilde e_i,V\rangle$. 度量相容及无挠性给出
\[
 \dd\eta_i(V,W)=\langle\nabla_V\widetilde e_i,W\rangle-
 \langle\nabla_W\widetilde e_i,V\rangle=0.
\]
再用同样的局部积分法得到函数 $y_i$ 使 $\dd y_i=\eta_i$. 两一形式独立, 逆函数定理使 $(y_1,y_2)$ 成为局部坐标. 任意切向量的平方长度为 $\eta_1(V)^2+\eta_2(V)^2$, 故度量为 $\dd y_1^2+\dd y_2^2$, 映射到欧氏平面是局部等距. 单连通小邻域避免闭形式的周期障碍; 没有由此声称整张平直曲面全局等距于平面.
\end{solution}
\par\Needspace{9\baselineskip}
\originalsection{测地线}
\par\Needspace{5\baselineskip}
\noindent\textbf{习题\ref{ex:8a}.}
\begin{solution}
圆柱路径的水平单位方向可由引理\ref{lem:angle}取连续角函数, 故可以沿整条路径取展开坐标. 展开覆盖把终点提升为 $(R(\theta+2\pi m),h)$, $m\in\Z$. 从起点 $(0,0)$ 到每个提升点的直线投影给出一条测地线, 其长度为
\[
 L_m=\sqrt{R^2(\theta+2\pi m)^2+h^2}.
\]
反之每条圆柱测地线提升后是直线, 所以这些列全. 长度最小要求 $|\theta+2\pi m|$ 最小. 当 $0\le\theta<\pi$, 取 $m=0$; 当 $\pi<\theta<2\pi$, 取 $m=-1$; 当 $\theta=\pi$, 两者相同, 恰有两条不同最短轨迹. $\theta=0,h=0$ 为同点情形, 最短是常值路径, 非平凡绕圈都较长. 任意曲面路径提升后的欧氏长度不小于相应端点距离, 因而该比较也排除了非测地竞争路径.
\end{solution}
\par\Needspace{5\baselineskip}
\noindent\textbf{习题\ref{ex:8b}.}
\begin{solution}
由 $\dot\theta=c/r^2$, $\dot s=\pm\sqrt{1-c^2/r^2}$, 在 $\dot s\ne0$ 的区间有
\[
 \frac{\dd\theta}{\dd s}=\frac{\pm c}{r\sqrt{r^2-c^2}},
\]
符号与 $\dot s$ 相同. 在 $r=|c|$ 只知道经向速度瞬时为0, 可能是转向点. 要使整个纬线 $s=s_0$ 成为测地线, 必须由经向方程得到 $r'(s_0)=0$; 此条件也充分. 若 $r'\ne0$, $\ddot s=rr'\dot\theta^2$ 非零, 路径立即离开该纬线.
\end{solution}
\par\Needspace{5\baselineskip}
\noindent\textbf{习题\ref{ex:8c}.}
\begin{solution}
记反射限制为等距 $\sigma:M\to M$. 选固定曲线一点 $p$ 和单位切向 $V$. 因 $\sigma$ 在该曲线上恒等, $\sigma(p)=p,\dd\sigma_pV=V$. 具有初值 $p,V$ 的测地线 $\gamma$ 经等距后仍为测地线; 这来自能量第一变分或度量联络保持. $\sigma\circ\gamma$ 有相同初值, 唯一性给出 $\sigma\circ\gamma=\gamma$. 它因此留在局部固定点集.

题设固定点集局部是一条正则曲线, 所以此非恒定测地线在短时间内正好沿该曲线行进. 采用相同方向弧长参数的固定曲线与它相同, 因而为测地线. 对整曲线逐点应用即可. 正则固定曲线条件排除固定点集只有一点、具有分叉或为整个二维曲面的情形; 后一种情形不能据此把任意曲线都判为测地线.
\end{solution}
\par\Needspace{9\baselineskip}
\originalsection{距离与正常坐标}
\par\Needspace{5\baselineskip}
\noindent\textbf{习题\ref{ex:9a}.}
\begin{solution}
球面为闭连通嵌入曲面, 定理\ref{thm:closed-geodesic}给出最短测地线. 前章已求出所有非恒定球面测地线为大圆. 长度为 $\ell$ 的单位速度段满足 $\langle p,q\rangle/R^2=\cos(\ell/R)=\cos\alpha$. 非负长度中最小解为 $R\alpha$, 而大圆较短弧确实达到, 所以距离为此值.

若 $0<\alpha<\pi$, 两点与球心确定唯一过心平面, 较短大圆弧轨迹唯一. 若 $\alpha=\pi$, 任一包含直径 $p,-p$ 的过心平面给出两个半圆弧, 总共有无穷多条最短轨迹. 若 $\alpha=0$, 最短长度为0, 只有常值路径; 非恒定大圆绕圈不最短. 这些唯一性是轨迹及单调参数意义, 同一轨迹允许不同恒速区间表示.
\end{solution}
\par\Needspace{5\baselineskip}
\noindent\textbf{习题\ref{ex:9b}.}
\begin{solution}
任意路径长度至少为位移长度2. 对 $0<\eps<1$, 先从 $(-1,0)$ 沿实轴到 $(-\eps,0)$, 再走上半圆半径 $\eps$, 最后从 $(\eps,0)$ 走到 $(1,0)$, 全部留在穿孔平面. 总长为 $2(1-\eps)+\pi\eps\to2$, 故下确界为2.

若有路径达到2, 欧氏长度不等式的取等要求其速度几乎处处为非负水平向量; 对分段光滑路径可逐段得到纵坐标恒为0且横坐标单调. 连续性迫使它经过原点, 与定义域矛盾. 即使允许一般可求长路径, 长度取等也使每个中间点落在线段上并按顺序前进, 同样经过原点.
\end{solution}
\par\Needspace{5\baselineskip}
\noindent\textbf{习题\ref{ex:9c}.}
\begin{solution}
光滑指数映射满足 $F(r,\theta)=p+r e(\theta)+O(r^2)$, 且可对 $\theta$ 微分, 故 $F_\theta=r e'(\theta)+O(r^2)$, $|e'|=1$. 所以 $A(0,\theta)=0$, $A_r(0,\theta)=1$; 更具体地, $F_\theta=rB(r,\theta)$, $B$ 光滑且 $B(0,\theta)=e'(\theta)\ne0$, 所以 $A=r|B|$ 在 $r\ge0$ 有单侧光滑延拓, 这些初值可用于积分微分方程.

对 $r>0$, 第7章的正交度量公式给出 $A_{rr}=-K(F(r,\theta))A$. 两次积分及初值得
\[
 A(r,\theta)=r-\int_0^r(r-t)K(F(t,\theta))A(t,\theta)\dd t.
\]
由初始展开 $A(t)=t+O(t^2)$ 和 $K$ 有界, 积分为 $O(r^3)$, 所以先得到 $A=r+O(r^3)$. 再用 $K(F(t,\theta))=K(p)+o(1)$, 有
\[
 A=r-K(p)\int_0^r(r-t)t\dd t+o(r^3)
  =r-\tfrac16K(p)r^3+o(r^3).
\]
当 $K>0$ 时, 小圆的环向尺度比欧氏小, 与球面 $A=R\sin(r/R)$ 一致; 负曲率时则比欧氏大.
\end{solution}
\par\Needspace{9\baselineskip}
\originalsection{平行移动与局部公式}
\par\Needspace{5\baselineskip}
\noindent\textbf{习题\ref{ex:10a}.}
\begin{solution}
保持曲线参数方向, $N\mapsto-N$ 使 $N\times T$ 反向, 因而 $k_g\mapsto-k_g$; $S,\II$ 变号使 $k_n\mapsto-k_n$; 两主曲率一起变号, 所以 $K$ 不变.

Gauss--Bonnet 的边界不是任取方向, 而是区域在定向切平面的左侧. 反转曲面定向后, 正向边界也要反向. 曲线反向使 $k_g$ 再变一次号, 两次相消, 所以按照新正向边界计算的测地曲率积分数值保持. 曲面内角不变, 以新曲面定向及新边界方向同时定义的外角也不变. Gauss 曲率乘普通正面积的积分保持, 公式因此一致.
\end{solution}
\par\Needspace{5\baselineskip}
\noindent\textbf{习题\ref{ex:10b}.}
\begin{solution}
过心平面的交线给出互相垂直的球面大圆, 三个顶点均为直角, 内角和 $3\pi/2$. 球面被这三平面分成八个全等区域, 三角形面积为 $4\pi R^2/8=\pi R^2/2$. 乘 $K=1/R^2$ 得曲率积分 $\pi/2$. 因而 $\pi+\int K=3\pi/2$, 等于角和. 三个外角均为 $\pi/2$, 其总和与曲率积分合起来为 $2\pi$.
\end{solution}
\par\Needspace{5\baselineskip}
\noindent\textbf{习题\ref{ex:10c}.}
\begin{solution}
帽区域面积为 $2\pi R^2(1-\cos\theta_0)$, 由 holonomy 公式, 净旋转角为 $2\pi(1-\cos\theta_0)$ 模 $2\pi$. 当 $\theta_0=\pi/2$, 实数积分为 $2\pi$, 净旋转在单位向量上等于0角; 一个整周转动的终点向量与初始向量相同. 非零曲率并不保证每一个回路都有非恒等 holonomy, 它保证足够小且有非零曲率积分的回路出现相应小转角.
\end{solution}
\par\Needspace{9\baselineskip}
\originalsection{整体公式}
\par\Needspace{5\baselineskip}
\noindent\textbf{习题\ref{ex:11a}.}
\begin{solution}
对连通曲面, $-\Area=\int K=4\pi(1-g)$, 故 $\Area=4\pi(g-1)$. 面积为正迫使 $g\ge2$. 这只是内蕴几何的必要拓扑与面积条件, 不自动产生三维欧氏嵌入; 事实上若它是紧致无边界嵌入曲面, 命题\ref{prop:elliptic-point}要求存在 $K>0$ 点, 与 $K=-1$ 矛盾. 抽象曲面与嵌入曲面的范围必须区分.
\end{solution}
\par\Needspace{5\baselineskip}
\noindent\textbf{习题\ref{ex:11b}.}
\begin{solution}
曲率积分为
\[
 \int K\dd A=\int_0^{2\pi}\int_{\theta_1}^{\theta_2}\sin\theta\dd\theta\dd\varphi
 =2\pi(\cos\theta_1-\cos\theta_2).
\]
在北边界 $\theta_1$, 区域位于更大的极角一侧, 正向边界沿 $\varphi$ 减小, 所以该边积分为 $-2\pi\cos\theta_1$. 南边界 $\theta_2$ 沿 $\varphi$ 增加, 积分为 $2\pi\cos\theta_2$. 两者之和恰抵消曲率积分. 纬带为环形区域, $\chi=0$, 无角点, 故完全满足整体公式.
\end{solution}
\par\Needspace{5\baselineskip}
\noindent\textbf{习题\ref{ex:11c}.}
\begin{solution}
平面曲率 $K=0$, 环形区域 $\chi=0$, 无角点, 所以两边界测地曲率积分总和为0. 若取向上法向, 外边界使区域在左, 按逆时针走, 是正向简单闭曲线, 转角定理给 $2\pi$. 内边界为使环形区域在左必须顺时针走, 积分为 $-2\pi$. 同一条内边界若单独作为其所围小圆盘的边界会取相反方向; 不能忽略其所界区域而只凭曲线形状定向.
\end{solution}
\par\Needspace{9\baselineskip}
\originalsection{双曲平面}
\par\Needspace{5\baselineskip}
\noindent\textbf{习题\ref{ex:12a}.}
\begin{solution}
由距离公式, $d(0,1/2)=2\arctanh(1/2)=\log3$. 两点 $-1/2,1/2$ 对应的 $q=1/(1+1/4)=4/5$, 所以 $d(-1/2,1/2)=\log9=2\log3$. 沿实轴, $x(s)=\tanh((s-s_0)/2)$ 有 $2|x'|/(1-x^2)=1$, 是单位速度参数, 可选 $s_0$ 调整起点. 方向反向取其负号或反向参数.
\end{solution}
\par\Needspace{5\baselineskip}
\noindent\textbf{习题\ref{ex:12b}.}
\begin{solution}
三角形角和为 $\pi(1/3+1/4+1/6)=3\pi/4$, 双曲面积为 $\pi/4$. 若单位球面中有这样的测地圆盘三角形, 则 $K=1$ 的 Gauss--Bonnet 会给面积等于角和减 $\pi=-\pi/4$, 不可能. 所以不存在所指球面三角形. 此判断按简单测地圆盘和实际内角作出, 不把跨多周或带重数的三条弧称为同一种三角形.
\end{solution}
\par\Needspace{5\baselineskip}
\noindent\textbf{习题\ref{ex:12c}.}
\begin{solution}
度量是 $e^{2f}(\dd x^2+\dd y^2)$, $f=-\log y$, $\Delta f=1/y^2$, 所以 $K=-1$. Christoffel 公式给出 $\Gamma^x_{xy}=\Gamma^x_{yx}=-1/y$, $\Gamma^y_{xx}=1/y$, $\Gamma^y_{yy}=-1/y$, 其余为0. 若 $x=x_0$, 测地线方程变为 $y''-y'^2/y=0$, 解 $y(t)=Ae^{ct}$, $A>0$. 速度为 $|c|$, 所以 $c=\pm1$ 是单位速度.

任意连接同一竖线两点的路径长度满足
\[
 L=\int\frac{\sqrt{x'^2+y'^2}}y\dd t
 \ge\int\frac{|y'|}y\dd t
 \ge|\log y_2-\log y_1|.
\]
竖直单调路径取等, 因此距离为 $|\log(y_2/y_1)|$. 第一不等式取等要求 $x'=0$, 第二要求高度不回走, 也给出最短轨迹唯一.
\end{solution}

\originalchapter{官方讲义对应与范围}
官方四份讲义共60个物理页, 含每份的 MIT OCW 来源页及章标题页, 数学正文页数不同于课程讲次. 下列对应基于实际读到的讲义, 不由课程名称推测. 第27--29讲的逆函数和水平集内容在原件有重复, 本书合并到曲面定义章.
\begin{longtable}{p{.18\textwidth}p{.39\textwidth}p{.34\textwidth}}
\toprule
官方讲义与讲次 & 原件实际内容 & 本书处理\\
\midrule\endhead
第1章, 1--3 & 平面曲率、弧长、Frenet、曲线基本定理、密切圆、水平集曲率 & 第1章补条件与证明, 配椭圆和隐式算例.\\
第1章, 4 & Minkowski 平面曲线 & 不纳入欧氏曲线主线.\\
第1章, 5--6 & 高维 Frenet 标架及三维曲率、挠率 & 第3章作三维特化, 补曲线基本定理与螺线判据.\\
第1章, 7--10 & 圆映射次数、绕数、旋转数、Jordan、转角、凸性、Whitney、四顶点 & 第2章取切角、转角、严格正曲率凸性和四顶点; 不给 Whitney 分类或一般正则同伦理论.\\
第2章, 11--15 & 正定内积、超曲面基本形式、形算子、主曲率、旋转及切线曲面 & 第4--6章作二维曲面特化, 明示平均曲率因子差异.\\
第2章, 16--18,21 & Christoffel、Gauss 方程、曲率内蕴性、移动标架、特殊坐标 & 第7章逐式证明, 增补 Codazzi 和曲面基本定理局部存在性. 正常坐标在第9章构造.\\
第2章, 19--20,22--23 & 高维外代数曲率、刚性、Minkowski、余维推广、抽象度量 & 高维与 Lorentz 推广不纳入; 第12章使用二维圆盘内蕴度量.\\
第3章, 24--30 & 水平集、切空间、定向、图像化、紧致超曲面、Hadamard & 第4--5、11章二维特化; 逆隐函数为先修, 凸性结论给几何证明并明示覆盖工具.\\
第3章, 31--35 & 面积、映射次数、Gauss 映射总曲率、带奇点标架、Gauss--Bonnet、组合曲率 & 第4、11章取曲面面积和 Gauss Jacobian. 用局部公式和三角剖分证明整体 Gauss--Bonnet. 不系统讲一般次数、Hopf 高维版本或多面体重复内容.\\
第4章, 36--39 & 测地线、旋转曲面、变分、Hamilton、内蕴距离、双曲距离、Schwarz--Pick & 第8--9章证明测地线与最短性. 第12章证明圆盘距离. 不扩展 Hamilton 动力系统或复分析 Schwarz--Pick.\\
第4章, 40--41 & Busemann/CAT、测地曲率、带边界 Gauss--Bonnet、测地三角形 & 第10--11章完整几何公式及角点; 度量曲率空间不纳入.\\
\bottomrule
\end{longtable}
本册增加的教学材料包括图像曲率完整计算、悬链面和螺旋面、法向面积变分、Clairaut 关系细节、Gauss 引理、闭嵌入曲面的最短线存在证明以及36道练习的完整解答. 这些内容是对选定主线的重构与推导, 不伪称每一题均为原课程作业或原作者的逐字陈述. 本册不讨论一般流形的系统拓扑、向量丛特征类、高维 Riemann 几何、完整负曲率嵌入障碍、Jacobi 场与完整二阶变分理论; 这一区域的缺失不以“流形入门”替代.

\par\Needspace{8\baselineskip}
\originalsection{官方作业的覆盖与缺口}
10份公开Homework每份2个物理页, 共20页, 含10个题面页. 逐题实际计数为3、2、3、5、3、4、3、2、4、3, 共32个顶层题. 原件四题带有10个正式标号子问, 属于上述32题内部. 作业1--10按原次序附在基础正文与36题解答之后, 并保留原分值、条件和指涉. 30个顶层题的全部可恢复要求得到独立解答, 其中原题错误的两个命题给出反例与修正版. 作业3.3的折线定义、总曲率与Hopf部分已解, 但其最后所引Proposition 6.3没有可恢复陈述; 作业9.4仅引用未定位的Lemma 28.3, 因而保留缺口, 不猜题.

作业6.1的商业教材方法未复制, 改为明确署作独立推导的自足翻转隆起构造. 作业7.2的基本形式用词及原点归一化条件、作业10.1的讲次编号、作业10.2的有符号角动量不等式均保留原文并批注. 新增作业还补齐所用外幂、球面次数积分、奇点标架及受约束粒子方程; 这不等于系统展开整个高维课程. 公开目录未提供官方答案或考试原卷, 本册全部新增解答为AI独立编写.

\begin{thebibliography}{9}
\bibitem{seidel} Paul Seidel. \textit{18.950 Differential Geometry, Fall 2008}. MIT OpenCourseWare. \href{https://ocw.mit.edu/courses/18-950-differential-geometry-fall-2008/pages/lecture-notes/}{官方讲义目录}, 四章, 第1--41讲. 原课程参考书在官方 syllabus 登记, 本书未据未读教材编造页码或题源.
\bibitem{terms} MIT OpenCourseWare. \href{https://ocw.mit.edu/pages/privacy-and-terms-of-use/}{Privacy and Terms of Use}. 核对日期2026-10-08. \href{https://creativecommons.org/licenses/by-nc-sa/4.0/}{CC BY-NC-SA 4.0 International}.
\end{thebibliography}
\end{document}
